\chapter{多变量函数的微分学}
\section{多变量函数及其连续性}


\subsection{点集的结构}
多变量函数是定义在高维空间点集上的函数，为了满足所要研究函数的性质，需要对这个点集所具有的结构做进一步的规定.

\begin{definition}
    定义集合$\mathbb{R}^n$是有序$n$元实数组的集合：
    \[\mathbb{R}^n \coloneq \{(x_1,x_2,\cdots,x_n)\,|\,x_i \in \mathbb{R}\}\]
    其中的元素$\vb{x} = (x_1,x_2,\cdots,x_n)$称为$\mathbb{R}^n$中的一个点.
\end{definition}

\subsubsection{实线性空间}
\begin{definition}
    对于$\mathbb{R}^n$中的两点$\vb{x}=(x_1,\cdots,x_n),\vb{y}=(y_1,\cdots,y_n)$，定义加法运算：
    \[\vb{x} + \vb{y} \coloneq (x_1+y_1,\cdots,x_n+y_n)\]
    定义数乘运算：
    \[\forall \lambda \in \mathbb{R},\quad \lambda\vb{x} \coloneq (\lambda x_1,\cdots,\lambda x_n)\]
    定义有上述运算的空间称为\overtext{实线性空间}{real linear space}或\overtext{实向量空间}{real vector space}.
\end{definition}


\subsubsection{赋范空间}
\begin{definition}
    在线性空间上定义\overtext{范数}{norm}运算：
    \[\left\|\vb{x}\right\| \coloneq \sqrt{x_1^2 + \cdots + x_n^2}\]
    称定义有范数的线性空间为\overtext{赋范空间}{normed space}\!\wiki{https://zh.wikipedia.org/wiki/\%E8\%B3\%A6\%E7\%AF\%84\%E5\%90\%91\%E9\%87\%8F\%E7\%A9\%BA\%E9\%96\%93}.
\end{definition}


\subsubsection{度量空间}
\begin{definition}
    在赋范空间上定义\overtext{距离}{metric}的概念：
    \[\rho(\vb{x},\vb{y}) \coloneq \left\|\vb{x} - \vb{y}\right\| = \sqrt{\sum_{i=1}^{n}(x_i-y_i)^2}\]
    其满足距离的四条公理：
    \begin{enumerate}[circ]
        \item \textbf{非负性：}$\rho(\vb{x,y}) \geqslant 0$.
        \item \textbf{不可区分性}$\rho(\vb{x,y})\Leftrightarrow \vb{x=y}$.
        \item \textbf{对称性：}$\rho(\vb{x,y}) = \rho(\vb{y,x})$.
        \item \textbf{三角不等式：}$\rho(\vb{x,y})\leqslant\rho(\vb{x,z}) + \rho(\vb{y,z})$.
    \end{enumerate}
    称定义有距离的空间为\overtext{度量空间}{metric space}\!\wiki{https://zh.wikipedia.org/wiki/\%E5\%BA\%A6\%E9\%87\%8F\%E7\%A9\%BA\%E9\%97\%B4}.
\end{definition}


\subsubsection{拓扑空间}
为了定义极限，要为集合赋予「一点附近」的概念.这样的空间称为\overtext{拓扑空间}{topological space}\!\wiki{https://zh.wikipedia.org/wiki/\%E6\%8B\%93\%E6\%89\%91\%E7\%A9\%BA\%E9\%97\%B4}.在拓扑空间中，可以定义以下概念：

\begin{definition}
    设有一点$\vb{x}_0$，称满足如下定义的集合为一个\overtext{开球}{open ball}：
    \[B(\vb{x}_0,r) = \{\vb{x}\,|\,\rho(\vb{x},\vb{x}_0) < r\}\]
    除去中心点$\vb{x}_0$，则得到\overtext{去心开球}{punctured open ball}的定义：
    \[B_-(\vb{x}_0,r) = \{\vb{x}\,|\,0 < \rho(\vb{x},\vb{x}_0) < r\}\]
\end{definition}

\begin{definition}
    邻域是包含某个开球的集合：
    \begin{itemize}
        \item 称$U$是$\vb{x}_0$的\overtext{邻域}{neighborhood}，当且仅当$\exists r>0$使得$B(\vb{x}_0,r)\subseteq U$.
        \item 称$U_-$是$\vb{x}_0$的\overtext{去心邻域}{punctured neighborhood}，当且仅当$\exists r>0$使得$B_-(\vb{x}_0,r)\subseteq U_-$.
    \end{itemize}
\end{definition}

从邻域的定义不难发现，邻域只要能够包含某个开球即可，并没有固定的形状，常用的是球形或方形.

\begin{definition}
    若对于一个点集$E$，$\exists R>0$使得$E \subset B(O,R)$，则称$E$是\textbf{有界的}，其中$O$是坐标原点.

    定义有界非空点集$E$的\overtext{直径}{diameter}为集合中任意两点距离的上确界，记作：
    \[\mathrm{diam}\, E = \sup\{\rho({\vb{x,y}})\,|\,\vb{x,y}\in E\}\]
\end{definition}

\begin{definition}
    设$E \subset \mathbb{R}^n$是一个点集，$\mathbb{R}^n$中不属于$E$的部分称为$E$的\overtext{余集}{complement}，记作$E^c$，显然有$(E^c)^c = E$.平面上的点按照其与$E$的关系分为以下三种：
    \begin{enumerate}[circ]
        \item 如果$\exists r>0$，使得$B(\vb{x}_0,r) \subset E$，则称$\vb{x}_0$是$E$的\overtext{内点}{interior point}.$E$中全体内点组成的集合称为$E$的\overtext{核}{interior}，记作$E^\circ$.
        \item 如果$\exists r>0$，使得$B(\vb{x}_0,r) \subset E^c$，则称$\vb{x}_0$是$E$的\overtext{外点}{exterior point}.
        \item 如果$\forall r>0$，使得$B(\vb{x}_0,r)$中既有$E$的点，又有$E^c$的点，则称$\vb{x}_0$为$E$的\overtext{边界点}{boundary point}.$E$的全体边界点的集合称为$E$的\overtext{边界}{boundary}，记作$\partial E$.
    \end{enumerate}
\end{definition}

显然按照边界点的定义，边界点既可以在$E$中，也可以不在$E$中.

\begin{definition}
    \begin{itemize}
        \item 对于$E$中的点，如果$\exists r>0$，使得$B(\vb{x}_0,r) \cap E = \vb{x}_0$，那么称$\vb{x}_0$是$E$的\overtext{孤立点}{}.
        \item 对于空间中一点$\vb{x}_0$，如果$\forall r > 0$，使得$B_-(\vb{x}_0,r)$中总是有$E$中的点，称$\vb{x}_0$为$E$的\overtext{聚点}{}.
    \end{itemize}
\end{definition}

根据上述定义，$E$的聚点不一定在$E$中，$E$的边界点要么是聚点，要么是孤立点.不孤立的$E$的边界点则一定是聚点，在$E$中的内点也一定是聚点.

\begin{definition}
    设$E \subset \mathbb{R}^n$，若$E = E^\circ$，即$E$中的点都是$E$的内点，则称$E$是一个\overtext{开集}{open set}；若$E$的余集$E^c$是开集，则称$E$是\overtext{闭集}{close set}.
\end{definition}

按照上述定义，在$\mathbb{R}^n$中，只有整个空间$\mathbb{R}^n$和空集$\emptyset$既是\overtext{既开又闭集}{clopen set}.一方面，它们中没有不是内点的点，另一方面，二者互为余集，因此也是闭集.

\begin{definition}
    设$E \in \mathbb{R}^n$，记$\overline{E} = E \cup \{E\text{的所有聚点}\}$，称为$E$的\overtext{闭包}{}.
\end{definition}
可以认为$E$的闭包就是$E$本身加上所有从极限上能够靠近到的点，或者说是$E$中所有点加上不在$E$中的$E$的边界点.

开集和闭集具有如下性质：
\begin{enumerate}[property]
    \item 两个开集的并或交仍是开集，两个闭集的并或交仍是闭集.
    \begin{proof}
        以开集之交的证明为例：设有开集$E_1,E_2$，$\forall M \in E_1 \cap E_2$，满足$M \in E_1$且$M \in E_2$.对$M \in E_1$，$\exists r_1 > 0$使得$B(M,r_1) \subset E_1$；对$M \in E_2$，$\exists r_2 > 0$使得$B(M,r_2) \subset E_2$.

        不妨取$r = \min\{r_1,r_2\}$，则对$\forall M \in E_1 \cap E_2$，$\exists r$满足$B(M,r) \subset (E_1 \cap E_2)$.
    \end{proof}
    \item $E\subset\mathbb{R}^n$是开集，当且仅当$\partial E \cap E = \emptyset$.
    \item $E\subset\mathbb{R}^n$是闭集，当且仅当$\partial E \subset E$.
    \item $E\subset\mathbb{R}^n$是闭集，当且仅当$E = \overline{E}$.
\end{enumerate}
\begin{mistake}
    设有命题$p,q,q'$，若有$p\Leftrightarrow q$且$p\Leftrightarrow q'$，则命题$q$和$q'$在表述上可以不一致，但是在逻辑上一定等价，即$q \Leftrightarrow q'$.
\end{mistake}

\begin{definition}
    对于点集$E \subset \mathbb{R}^n$，如果存在两个非空开集$A,B$使得$E = A \cup B$时，则称$E$是\textbf{不}\!\!\overtext{连通}{connected}的，否则称$E$是连通的.
\end{definition}


\begin{definition}
    \textbf{连续曲线}是由如下参数方程表示函数的图像：
    \[\vb{r}(t) = (x(t),t(t)),\quad t\in[\alpha,\beta]\]
    其中$x(t),y(t)$均为$[\alpha,\beta]$上的连续函数.若上式满足对任意的$\alpha \leqslant t_1 < t_2 \leqslant \beta$有$(x(t_1),y(t_1)) \neq (x(t_2),y(t_2))$，即该曲线不存在自相交，则称曲线为一条\textbf{简单曲线(Jordan曲线)}.

    若曲线的两端点重合，即$(x(\alpha),y(\alpha)) = (x(\beta),y(\beta))$，则称为\textbf{闭曲线}.
\end{definition}

\begin{definition}
    对于点集$E$中任意两点$P,Q$，若在$E$中存在一条连续曲线将$P$和$Q$连接起来，则称$E$是\overtext{道路连通}{path connected}或\overtext{弧连通}{arc connected}的.
\end{definition}

\begin{theorem}
    道理连通集一定是连通的，反之则不然.
\end{theorem}

\begin{proof}
    取$E = A \cup B$且$A \cap B = \emptyset$.取$\forall M_1 \in A,\forall M_2 \in B$，因为$E$是道路连通的，因此存在连续曲线：
    \[\vb{r}(t) = (x(t),y(t)),\quad t\in[\alpha,\beta]\quad\text{且}\quad\vb{r}(\alpha) = M_1,\vb{r}(\beta) = M_2\]

    设$t_0 = \inf\{t\,|\,\vb{r}(t) \in B\}$，此时分两种情况讨论.若$\vb{r}(t_0)\in A$，则$\forall\varepsilon > 0$，$\exists\delta>0$使得在$|t-t_0|<\delta$时有：
    \[|x(t) - x(t_0)| < \frac{\varepsilon}{2},\qquad |y(t) - y(t_0)| < \frac{\varepsilon}{2}\]
    若是取$t=t',\quad t_0 < t < t_0 + \delta$，则显然$\vb{r}(t') \in B$.

    即在$B(\vb{r}(t_0),\varepsilon)$中对$\forall \varepsilon>0$都存在$B$中的点，则表明$\vb{r}(t_0)\in A$是$B$的聚点，满足连通的定义.
\end{proof}

\begin{definition}
    道路连通的开集称为\overtext{开区域}{open domain}，开区域的闭包称为\overtext{闭区域}{close region}，统称为\overtext{区域}{region}.
\end{definition}

\begin{definition}
    如果区域$D$内任意一条简单闭曲线的内部仍在$D$内，即$D$内没有空洞，那么称$D$是\overtext{单连通}{simply connected}的，反之则称为\overtext{多连通}{multiply connected}的.
\end{definition}


\subsection{平面上的点集}
\subsubsection{平面点列的极限}
\begin{definition}
    设$\{M_n\}$是平面点列，如有$M_0 \in \mathbb{R}^2$使得$\lim \rho(M_n,M_0) = 0$，则称点列$\{M_n\}$为\textbf{收敛点列}，$M_0$称为点列的\textbf{极限}.记作：
    \[\lim_{n \to \infty} M_n  = M_0\]
\end{definition}

平面上的点列不一定需要以直线逼近其极限，可以由任意的轨迹逼近.

\begin{theorem}
    设$M_n(x_n,y_n)$，则有：
    \[\lim M_n = M_0 \quad\Leftrightarrow\quad \begin{cases}
        \lim x_n = x_0 \\
        \lim y_n = y_0
    \end{cases}\]
\end{theorem}

\begin{proof}
    注意到有以下不等式：
    \begin{gather*}
        |x_n-x_0| \leqslant \rho(M_n,M_0) \leqslant |x_n-x_0| + |y_n-y_0| \\
        |y_n-y_0| \leqslant \rho(M_n,M_0) \leqslant |x_n-x_0| + |y_n-y_0|
    \end{gather*}
    根据三明治法则，即得证.
\end{proof}

\begin{theorem}
    \label{thm:2DBolzanoWeierstrass}
    平面上的有界点列必有收敛子列.
\end{theorem}

\begin{proof}
    设有$M_n(x_n,y_n)$，若$\{x_n\}$有界，则必有收敛子列$\{x_{n_k}\}$，则$\{x_{n_k}\}$的子列$\{x_{n_k'}\}$也收敛.

    若$\{y_n\}$有界，则必有子列$\{y_{n_k}\}$有界，则其进一步有有界子列$\{y_{n_k'}\}$收敛.

    则平面上点列$\{M_n\}$必有收敛子列$\{M_{n_k'}\}$.
\end{proof}


\subsection{多变量函数}
\subsubsection{多变量函数的定义}
$n$维实数组空间$\mathbb{R}^n = \{\vb{x} = (x_1,x_2,\cdots,x_n)\,|\,x_i\in\mathbb{R},i=1,2,\cdots,n\}$中的点$\vb{x}$对应一个向量：
\[\vb{x} = x_1\vb{e}_1 + x_2\vb{e}_2 + \cdots + x_n\vb{e}_n\]
其中$\vb{e}_i$是$\mathbb{R}_n$的标准正交基.

\begin{definition}
    设$D \subset \mathbb{R}^n$，则从$D$到直线上的映射称为一个\overtext{$n$元函数}{$n$-variable function}，记作：
    \[f:D \to \mathbb{R} \quad\text{或}\quad z = f(x_1,x_2,\cdots,x_n)\]
\end{definition}

对于二元函数$z=f(x,y),(x,y)\in D \subset \mathbb{R}^n$，如果把$z$所在的数轴垂直于$Oxy$平面放置并形成右手系，则得到一个三维空间$Oxyz$上的一张曲面：
\[S = \{(x,y,f(x,y))\,|\,(x,y)\in D\}\]
若在$Oxy$平面上取使得$f(x,y)$等于一固定常数$c$的点的集合：
\[L_c = \{(x,y)\,|\,f(x,y)=c\}\]
则得到了$Oxy$平面上的一条或几条曲线，或若干点，或为空集，称为\overtext{等值线}{contour line}.


\subsection{多变量函数的复合}
将两个二元函数复合，得到的仍是二元函数；将二元函数和一元函数复合得到的则是一元函数.
\begin{figure}[H]
    \centering
    \includegraphics[width=\textwidth]{composition-of-multivar-func/fig.pdf}
    \caption{多变量函数的复合映射过程}
\end{figure}


\subsection{向量值函数}
向量值函数就是一个从$n$维空间$\mathbb{R}^n$到$m$维空间$\mathbb{R}^m$的映射：
\[\vb{f}: D \to \mathbb{R}^m,\quad D \subset \mathbb{R}^n \quad\text{或}\quad \vb{f}:(x_1,x_2,\cdots,x_n)\mapsto(y_1,y_2,\cdots,y_m)\]
应当注意的是，向量值函数$\vb{f}$的每一个分量$y_i$都是所有自变量的多元函数：
\[y_i = f_i(x_1,x_2,\cdots,x_n),\quad i = 1,2,\cdots,m\]
当维数$m=n$时，如果映射可逆，那么称为相同维数空间之间的变换，可逆意味着变换过程中没有发生信息丢失.

常用的向量值函数有$\mathbb{R}\to\mathbb{R}^3$的映射：
\[\vb{r}(t) = x(t)\vb{i} + y(t)\vb{j} + z(t)\vb{k},\quad t\in[\alpha,\beta]\subset\mathbb{R}\]
这将一维的参数$t$映射到了三维空间的一条曲线上，即空间曲线的参数方程表示，称为\overtext{参数曲线}{parametric curve}.

还有$\mathbb{R}^2 \to \mathbb{R}^3$的映射：
\[\vb{r}(u,v) = x(u,v)\vb{i} + y(u,v)\vb{j} + z(u,v)\vb{k},\quad (u,v)\in D \subset \mathbb{R}^2\]
这是将二维参数数组$(u,v)$映射到了三维空间中的一张曲面上，称为\overtext{参数曲面}{parameter surface}.


\subsection{多变量函数的极限}
\begin{definition}
    设$z=f(x,y),(x,y)\in D$，$M_0(x_0,y_0)$是$D$的聚点.若是对$\forall \varepsilon > 0$，$\exists \delta > 0$，当$M(x,y)\in D$满足：
    \begin{gather*}
        \text{圆心去心邻域：}0 < \rho(M,M_0) = \sqrt{(x-x_0)^2 + (y-y_0)^2} < \delta\\
        \text{方形去心邻域：}|x-x_0|<\delta,|y-y_0|<\delta,\quad(x,y)\neq (x_0,y_0)
    \end{gather*}
    使得$|f(M)-a|<\varepsilon$，则称当$M$趋于$M_0$时，$f(M)$以$a$为\textbf{极限}，记作：
    \[\lim_{M \to M_0}f(M) = a \quad\text{或}\quad \lim_{(x,y)\to(x_0,y_0)}f(x,y) = a \quad\text{或}\quad \lim_{\substack{x \to x_0 \\ y \to y_0}}f(x,y) = a\]
\end{definition}

\begin{question}
    证明：$\displaystyle \lim_{(x,y)\to(x_0,y_0)}\frac{x^2y^2}{x^4+y^2} = 0$.
\end{question}

\begin{proof}
    \[\left|\frac{x^2y^2}{x^4+y^2}\right| = \frac{|y|}{|x^4+y^2|}\cdot|x^2 y| \leqslant \frac{|y|}{|x^4+y^2|}\cdot\frac{|x^4+y^2|}{2} = \frac{|y|}{2}\]
    若有极限成立，则上式应满足$\forall \varepsilon > 0$，$\exists\delta>0$使得$|x-0|<\delta,|y-0|<\delta$时有$\frac{|y|}{2} < \varepsilon$.故而取$\delta = 2\varepsilon$，使得上述极限存在且等于$0$.
\end{proof}

\begin{question}
    证明：$\displaystyle \lim_{(x,y)\to(x_0,y_0)}\frac{x^3+y^3}{x^2+y^2} = 0$.
\end{question}

\begin{proof}
    取$x = r\cos\theta, y = r\sin\theta, r \to 0$，有：
    \[\left|\frac{r^3(\cos^3\theta + \sin^3\theta)}{r^2}\right| \leqslant r\]
    对$\forall \varepsilon > 0$，取$\delta = \varepsilon$，使得$r = \rho((x,y),0) < \delta$时有：
    \[\left|\frac{x^3+y^3}{x^2+y^2} - 0\right| \leqslant \varepsilon\]
    故而极限存在且为$0$.
\end{proof}

\begin{theorem}
    设$z=f(x,y),(x,y)\in D$，有$D(x_0,y_0)$是$D$的一个聚点，则有函数极限：
    \[\lim_{(x,y)\to(x_0,y_0)}f(x,y) = a\]
    的充分必要条件是对任何收敛于$(x_0,y_0)$的点列$(x_n,y_n),(x_n,y_n)\neq(x_0,y_0)$都满足：
    \[\lim_{n \to \infty}f(x_n,y_n) = a\]
\end{theorem}

根据\cref{subsec:LimitProperty}中性质7的证明，分别对上述定理中的$x,y$分别证明，而后合并为$\rho((x,y),(x_0,y_0))$进一步证明即可.

\begin{theorem}
    设$z = f(u,v)$，其中$u=u(x,y),v=v(x,y)$，若：
    \begin{enumerate}[circ]
        \item $\lim\limits_{(u,v) \to (u_0,v_0)}f(u,v) = a$
        \vspace{0.5em}
        \item $\lim\limits_{(x,y) \to (x_0,y_0)}u(x,y) = u_0,\quad \lim\limits_{(x,y) \to (x_0,y_0)}v(x,y) = v_0$
    \end{enumerate}
    则$\lim\limits_{(x,y) \to (x_0,y_0)}f(u(x,y),v(x,y)) = a$.
\end{theorem}

\begin{proof}
    对于$\forall \varepsilon > 0$，根据条件\circnum{1}，必定$\exists\delta' > 0$，使得当$|u-u_0|<\delta',|v-v_0|<\delta'$时，有$|f(u,v)-a|<\varepsilon$.

    根据条件\circnum{2}，对于一个给定的$\delta'>0$，一定$\exists \delta>0$在$|x-x_0|<\delta,|y-y_0|<\delta$且$(x,y)\neq(x_0,y_0)$时，满足$|u-u_0|<\delta',|v-v_0|<\delta'$.

    综上所述，对于$\forall \varepsilon>0$，$\exists \delta>0$，在$|x-x_0|<\delta,|y-y_0|<\delta$且$(x,y)\neq(x_0,y_0)$时，满足：
    \[|f(u(x,y),v(x,y))-a|<\varepsilon \quad\Rightarrow\quad \lim_{(x,y) \to (x_0,y_0)}f(u(x,y),v(x,y)) = a\]
\end{proof}

\begin{definition}
    设函数$f(x,y)$在$(x_0,y_0)$附近有定义，先后对$x,y$求极限：
    \[\lim_{y \to y_0} \lim_{x \to x_0}f(x,y),\qquad \lim_{x \to x_0} \lim_{y \to y_0}f(x,y)\]
    如果上述极限存在，那么称之为函数$f(x,y)$在$(x_0,y_0)$处的\overtext{累次极限}{iterated limit}.
\end{definition}

通过举反例不难说明，累次极限与多变量函数的极限没有必然的联系，一个存在不代表另一个也存在.

\begin{theorem}
    设有$\lim\limits_{(x,y) \to (x_0,y_0)}f(x,y) = A$，如果对$y\neq y_0$，有$\varphi(y) = \lim\limits_{x \to x_0}f(x,y)$存在，则：
    \[\lim_{y \to y_0}\left(\lim_{x \to x_0}f(x,y)\right) = A\]
\end{theorem}

\begin{proof}
    对$\forall \varepsilon>0$，$\exists \delta > 0$，在$|x-x_0|<\delta,|y-y_0|<\delta$时满足：
    \[|f(x,y)-A| < \frac{\varepsilon}{2} \quad\Leftrightarrow\quad A - \frac{\varepsilon}{2} < f(x,y) < A + \frac{\varepsilon}{2}\]
    令$x \to x_0$，上式变为：
    \[A - \frac{\varepsilon}{2} \leqslant \varphi(y) \leqslant A - \frac{\varepsilon}{2} \quad\Rightarrow\quad |\varphi(y) - A| \leqslant \frac{\varepsilon}{2} < \varepsilon\]
    即对$\forall \varepsilon>0$，$\exists\delta>0$在$|y-y_0|$时满足$|\varphi(y) - A| < \varepsilon$，即：
    \[\lim_{y \to y_0}\varphi(y) = A \quad\Rightarrow\quad \lim_{y \to y_0}\left(\lim_{x \to x_0}f(x,y)\right) = A\]
    由此得证.
\end{proof}

\begin{theorem}
    设$f(x,y)$在$M_0(x_0,y_0)$附近$B_-(M_0,r)$有定义，且在$M_0$处有有限极限$a$，$x=x(t),y=y(t)$在$0 < |t-t_0| < l$有定义，使得$\{(x(t),y(t))\,|\,0<|t-t_0|<l\}\subset B_-(M_0,r)$且$\lim\limits_{t \to t_0}x(t) = x_0,\lim\limits_{t \to t_0}y(t) = y_0$，则：
    \[\lim_{t \to t_0}f(x(t),y(t)) = a\]
\end{theorem}

\begin{proof}
    已知：
    \[\lim_{(x,y)\to(x_0,y_0)}f(x,y) = a\]
    即$\forall \varepsilon>0$，$\exists \delta_1 > 0$在$|\rho((x,y),(x_0,y_0))|<\delta_1$时满足$|f(x,y)-a|<\varepsilon$.

    现选取一条任意路径$x=x(t),y=y(t)$，在$t\to t_0$时有$x \to x_0,y \to y_0$.即对$\forall \varepsilon > 0$，$\exists \delta_2>0$，使得$0 < \sqrt{(x(t)-x_0)^2+(y(t)-y_0)^2} = \rho((x,y),(x_0,y_0)) < \delta_2$时满足：
    \[|x(t)-x_0|<\varepsilon,\qquad |y(t)-y_0|<\varepsilon\]
    不妨取$\delta = \min\{\delta_1,\delta_2\}$，即这个$\delta$对任意一条路径在$|\rho((x,y),(x_0,y_0))|<\delta$时满足：
    \[\lim_{t \to t_0}f(x(t),y(t)) = a\]
\end{proof}

上述定理表明，只要极限存在，则平面上的点沿任何路径趋于$(x_0,y_0)$时，极限都存在且一致.这个定理主要用于证明极限不存在的情形，只要能找到一条路径使得极限不存在或不一致即可.


\subsection{多变量函数的连续性}
\subsubsection{多元函数连续性与一致连续性的定义}
\begin{definition}
    设$f(x,y)$在$M_0(x_0,y_0)$的邻域$B(M_0,r)$上有定义，若存在极限：
    \[\lim_{\substack{x \to x_0 \\ y \to y_0}}f(x,y) = f(x_0,y_0) \quad\text{或}\quad \lim_{\substack{\Delta x \to 0 \\ \Delta y \to 0}}f(x_0+\Delta x, y_0+\Delta y) = f(x_0,y_0)\]
    则称$f$在$(x_0,y_0)$处\textbf{连续}.
\end{definition}

\begin{question}
    证明在定义域$xy>-1$内，下列函数在$(0,0)$处连续：
    \[f(x,y) = \begin{cases}
        \ln(1+xy)^{\frac{1}{xy}},\quad xy \neq 0\\
        1,\quad xy = 0
    \end{cases}\]
\end{question}

\begin{proof}
    注意到有：
    \[\lim\limits_{u \to 0}\ln(1+u)^{\frac{1}{u}} = 1\]
    即对$\forall \varepsilon > 0$，$\exists \delta'>0$，使得$|u| < \delta'$时满足：
    \[\left|\ln(1+u)^{\frac{1}{u}} - 1\right| < \varepsilon\]
    取$\delta = \sqrt{\delta'} > 0$，在$|x|<\delta,|y|<\delta$时有$|xy|<\delta'$，进而满足：
    \[|f(x,y)-f(0,0)| = \left|\ln(1+xy)^{\frac{1}{xy}} - 1\right| < \varepsilon\]
    故而函数在$(0,0)$处连续.
\end{proof}

从上述例题不难看出，如果多元函数在一点连续，在该点固定其它变量得到的单变量函数也在这一点连续，反之则不然.

例如对于函数：
\[f(x,y) = \begin{cases}
    \displaystyle \frac{2xy}{x^2+y^2},\quad (x,y) \neq (0,0)\\
    0,\quad (x,y) = (0,0)
\end{cases}\]
对于单变量函数$f(x,0),f(0,y)$都在$x=0$或$y=0$时连续，但是若是以另一条轨迹，如沿着$y=x$向$(0,0)$逼近，则函数不连续：
\[\lim_{\substack{x \to 0 \\ y \to 0}}\frac{2xy}{x^2+y^2} = 1 \neq f(0,0)\]
同样地，在多变量函数中也存在一致连续性的概念：
\begin{definition}
    设$f$在区域$D$上连续，对$\forall \varepsilon > 0$，$\exists \delta>0$使得区域上任意两点$(x,y)$和$(x',y')$在满足：
    \[\rho = \sqrt{(x-x')^2 + (y-y')^2} < \delta \quad\text{或}\quad |x'-x|<\delta,|y'-y|<\delta\]
    时使得$|f(x',y')-f(x,y)|<\varepsilon$，就称$f$在$D$上\textbf{一致连续}.
\end{definition}


\subsubsection{连续多元函数的性质}
\begin{enumerate}[property]
    \item 连续函数的加、减、乘、除(分母不为$0$)也是连续函数.\\
    特别地，对于除法：
    \[\lim\limits_{(x,y) \to (x_0,y_0)}\frac{f(x,y)}{g(x,y)} = \frac{f(x_0,y_0)}{g(x_0,y_0)}\]
    若是$f(x,y)$在$D$上连续，并且有$f(x_0,y_0) > 0$则一定存在$(x_0,y_0)$的一个邻域$I_{x_0}$，使得当$(x,y)\in I_{x_0}$时，$f(x,y) > 0$.这就像一块布，提起一点周围也会跟着升高，这就是连续函数极限的局部保号性.
    \item 连续函数的复合仍然连续.
    \item \textbf{(介值定理)}如果连续函数$f(M)$的定义域$D$是道路连通的，那么对于任意两点$M_1,M_2\in D$，$f$在$D$中能够取到$f(M_1)$和$f(M_2)$中的所有值.
    \begin{proof}
        由于$f(M)$的定义域$D$是道路连通的，故而必然存在连续曲线：
        \[L:\quad\vb{r}(t)=(x(t),y(t)),\quad t\in[\alpha,\beta]\quad\begin{cases}
            \vb{r}(\alpha) = M_1\\
            \vb{r}(\beta) = M_2
        \end{cases}\]
        对单变量函数$\vb{r}(t)$，作为连续函数的复合，其在$[\alpha,\beta]$上连续，因而对于$\forall c$满足$f(M_1) < c < f(M_2)$，根据\cref{thm:IntermediateValueTheorem}\textbf{(单变量连续函数的介值定理)}，总是有$t_0\in[\alpha,\beta]$使得$f(\vb{r}(t_0)) = c$：
        \[f(\vb{r}(\alpha)) = f(M_1) < f(\vb{r}(t_0)) < f(M_2) = f(\vb{r}(\beta))\]
    \end{proof}
    \item \textbf{(最值定理)}有界闭集上的连续函数一定能取到最大值或最小值.
    \begin{proof}
        假设连续函数$f(x,y)$在有界闭集$D$上无上界.则对于$\forall n$，总是$\exists (x_n,y_n) \in D$，使得$f(x_n,y_n) > n$.
            
        对于有界闭集上的点列$(x_n,y_n)$，根据\cref{thm:2DBolzanoWeierstrass}，其必有收敛子列$(x_{n_k},y_{n_k}) \to (x_0,y_0)$，且收敛点仍在$D$内.

        根据$f(x,y)$的连续性，有：
        \[f(x_0,y_0) = \lim\limits_{k \to \infty}f(x_{n_k},y_{n_k}) \geqslant \lim\limits_{k \to \infty}n_k = \infty\]
        上述结果与$f(x,y)$连续的定义相矛盾，故而连续的$f(x,y)$在其定义域$D$上必有上界.

        有上界则必有上确界，记：
        \[A = \sup_{(x,y)\in D}f(x,y)\]
        则对$\forall n>0$有$A-\frac{1}{n}$不是上确界，则$\exists (x_n,y_n) \in D$使得：
        \[A \geqslant f(x_n,y_n) > A-\frac{1}{n}\]
        则根据三明治定理$n \to \infty$时，$f(x,y)$能够取到最大值.最小值的证明同理.
    \end{proof}
    \cref{thm:TheMostValuableTheorem}是上述定理的单变量形式.
    \item \textbf{(Heine-Cantor定理)}有界闭集上的连续函数一定是一致连续的.
    \begin{proof}
        假设$f(x,y)$在有界闭集$D$上连续但不一致连续，那么$\exists\varepsilon_0 > 0$，取$\delta = \frac{1}{n}$，其中$n$为任意正整数，在$D$中不妨取两个点列$\{(x_n,y_n)\}$与$\{(x'_n,y'_n)\}$满足：
        \[|x'_n - x_n| < \delta = \frac{1}{n},\quad |y'_n-y_n| < \delta = \frac{1}{n}\]
        但是又有：
        \[|f(x'_n,y'_n) - f(x_n,y_n)| \geqslant \varepsilon_0\]
        由于有界闭集中的点列必有收敛子列，所以$\{(x_n,y_n)\}$与$\{(x'_n,y'_n)\}$可以各取出一个序号相同的收敛子列，并满足：
        \[|x'_{n_k} - x_{n_k}| < \delta = \frac{1}{n_k},\quad |y'_{n_k}-y_{n_k}| < \delta = \frac{1}{n_k}\]
        由三明治定理可知$k \to \infty$时两个子列收敛于同一点$(x_0,y_0)$，于是有：
        \[\lim\limits_{k \to \infty}|f(x'_{n_k},y'_{n_k}) - f(x_{n_k},y_{n_k})| = |f(x_0,y_0) - f(x_0,y_0)| = 0 \geqslant \varepsilon_0\]
        与$\varepsilon > 0$矛盾，故而$f$在$D$上是一致连续的.
    \end{proof}
    \cref{thm:HeineCantorTheorem}是上述定理的单变量形式.
\end{enumerate}


\section{多变量函数的微分}
\subsection{多变量函数的偏微商}
\begin{definition}
    设$f(x,y)$在$M_0(x_0,y_0)$的邻域中有定义，如果函数沿$x$方向变化率的极限存在，则称它为$z=f(x,y)$在$M_0$关于$x$的\overtext{偏导数}{partial derivative}或\textbf{偏微商}，记作：
    \[f'_x(x_0,y_0) = \pdv{f}{x} \coloneq \lim\limits_{x \to x_0}\frac{f(x,y_0) - f(x_0,y_0)}{x-x_0} = \lim\limits_{\Delta x \to 0}\frac{f(x_0+\Delta x,y_0)-f(x_0,y_0)}{\Delta x}\]
\end{definition}
如果在区域$D$中的每一点都有关于某个变量的偏导数，由此就在$D$上形成了一个二元函数，即偏导函数.

从几何上看，在二元函数所形成的曲面上用平行于$x$轴的平面截得一条曲线，则关于$x$的偏导数就是这条曲线上一点处的斜率.

\begin{mistake}
    在单变量函数中，可导必定连续.在多变量函数中，可以求偏导只表明在所求偏导的方向上连续，并不表示在所有方向上连续.
\end{mistake}

\begin{definition}
    设函数$f(x,y)$在定义域内每一点都有偏导数$\displaystyle\pdv{f}{x},\pdv{f}{y}$.如果它们仍有偏导数，则称得到的结果为\textbf{高阶偏导数}.

    对不同自变量分别求导的高阶偏导数称为\textbf{混合偏导数}，例如对二阶偏导数有：
    \[\pdv{x}\left(\pdv{f}{y}\right) = \pdv{f}{x}{y} = f''_{yx} = f''_{21},\quad \pdv{y}\left(\pdv{f}{x}\right) = \pdv{f}{y}{x} = f''_{xy} = f''_{12}\]
\end{definition}

混合偏导数求导的顺序不同，得到的结果不一定相同.

\begin{theorem}
    设$f(x,y)$在区域$D$内有定义，如果$\displaystyle\pdv{f}{x}{y}$和$\displaystyle\pdv{f}{y}{x}$连续，那么二者相等，或者说求导的次序可交换.
\end{theorem}

\begin{proof}
    取函数$f(x,y)$定义域$D$上一点$(x_0,y_0)$，设有$(x_0+\Delta x, y_0+\Delta y) \in B((x_0,y_0),r) \subset D$.

    令$f(x,y)$分别对$x,y$的一阶差分为：
    \[\begin{cases}
        \varphi(x) = f(x,y_0+\Delta y) - f(x,y_0)\\
        \psi(y) = f(x_0+\Delta,y) - f(x_0,y)
    \end{cases}\]
    注意到，二者分别再对另一个变量取差分后相等：
    \begin{align*}
        &\varphi(x_0+\Delta x) - \varphi(x_0) = \psi(y_0+\Delta y) - \psi(y_0)\\
        =& f(x_0+\Delta x, y_0+\Delta y) - f(x_0+\Delta x,y_0) - f(x_0,y_0+\Delta y) + f(x_0,y_0)
    \end{align*}
    根据微分中值定理有：
    \begin{align*}
        \varphi(x_0+\Delta x) - \varphi(x_0)
        =& \Delta x\varphi'(x_0+\theta_1\Delta x) \\
        =& \Delta x[f'_x(x_0+\theta_1\Delta x,y_0+\Delta y) - f'_x(x_0+\theta_1\Delta x,y_0)] \\
        =& \Delta x \Delta y f''_{xy}(x_0+\theta_1 \Delta x, y_0+\eta_1 \Delta y)
    \end{align*}
    其中$\theta_1,\eta_1 \in [0,1]$，类似地有：
    \[\psi(y_0+\Delta y) - \psi(y_0) = \Delta x \Delta y f''_{yx}(x_0 + \theta_2\Delta x,y_0+\eta_2\Delta y)\]
    于是令$\Delta x \to 0,\Delta y \to 0$，进而有：
    \begin{align*}
        f''_{xy}(x_0+\theta_1\Delta x,y_0+\eta_1\Delta y) &= f''_{yx}(x_0+\theta_2\Delta x, y_0+\eta_2\Delta y) \\
        f''_{xy}(x_0,y_0) &= f''_{yx}(x_0,y_0)
    \end{align*}
    由此得证.
\end{proof}


\subsection{多变量函数的可微性}
与单变量函数微分的定义类似，多变量函数的微分定义如下：
\begin{definition}
    设$f(x,y)$是定义在区域$D \subset \mathbb{R}^2$上的二元函数，如果在$D$中一点$(x_0,y_0)$的附近，函数$f(x,y)$所表示的曲面可以近似于过$f(x_0,y_0)$的一个平面，即有：
    \[f(x,y) = f(x_0,y_0) + a(x-x_0) + b(y-y_0) + o(\rho)\]
    其中$a,b$为常数，$\rho = \sqrt{(x-x_0)^2 + (y-y_0)^2}$.则称$f(x,y)$在$(x_0,y_0)$\textbf{可微}.
\end{definition}

设有二元函数作为一个映射：
\[f:\mathbb R^2\to\mathbb R,\qquad f=f(x,y)\]
并记自变量增量向量为：
\[\Delta\vb{\rho} \coloneq (\Delta x,\Delta y),\qquad \rho\coloneq\left\|\Delta\vb{\rho}\right\|\]
若存在唯一的线性映射$\mathcal A:\mathbb R^2\to\mathbb R$，使得：
\[f(x+\Delta x,y+\Delta y)-f(x,y) = \mathcal A(\Delta\vb{\rho})+o(\rho),\qquad\rho\to0\]
则称$f$在点$(x,y)$处可微，并定义微分映射
\[\dd: f\longmapsto \dd f\coloneq\mathcal A\]
因此，$\dd{f}$本身是一个从增量向量空间到函数值变化空间的线性映射，即：
\[\dd f:\mathbb R^2\to\mathbb R,\qquad\Delta f = \dd f(\Delta\vb{\rho})+o(\rho)\]
由于任意从$\mathbb R^2$到$\mathbb R$的线性映射均可表示为$\mathcal A(\Delta x,\Delta y) = a\Delta x+b\Delta y$，故有：
\[\dd f(\Delta\vb{\rho}) = a\Delta x+b\Delta y\]
另一方面，考虑两个坐标函数
\[\begin{cases}
    x:\mathbb R^2\to\mathbb R,\qquad x(x,y)=x\\
    y:\mathbb R^2\to\mathbb R,\qquad y(x,y)=y
\end{cases}\]
按照与$f$相同的定义，对坐标函数$x$有：
\[x(x+\Delta x,y+\Delta y)-x(x,y)=\Delta x\]
由于其增量本身就是关于$\Delta\vb{\rho}$的线性函数，因此有$\begin{cases}
    \dd x(\Delta\vb{\rho}) = \Delta x \\
    \dd y(\Delta\vb{\rho}) = \Delta y
\end{cases}$.于是，$\dd x,\dd y$就是两个基本的线性映射：
\[\dd x:\Delta\vb{\rho}\mapsto\Delta x,\qquad\dd y:\Delta\vb{\rho}\mapsto\Delta y\]
因此
\[\dd f(\Delta\vb{\rho})=a\dd{x}(\Delta\vb{\rho})+b\dd{y}(\Delta\vb{\rho}) \quad\Rightarrow\quad \dd{f} = a\dd{x} + b\dd{y}\]
在上述基于映射的定义的基础上，可以很好地将微分的定义推广到$n$元函数.

\begin{theorem}
    设函数$f(x,y)$在$(x_0,y_0)$可微，则有：
    \begin{enumerate}[circ]
        \item $f(x,y)$在$(x_0,y_0)$连续.
        \item $f(x,y)$在$(x_0,y_0)$的两个偏导数都存在，且$\displaystyle\pdv{f}{x}=a,\pdv{f}{y}=b$，即：
        \[\dd{f} = \pdv{f}{x}\dd{x} + \pdv{f}{y}\dd{y}\]
    \end{enumerate}
\end{theorem}

\begin{proof}
    因为$f$在$(x,y)$可微，则有：
    \[f(x+\Delta x,y+\Delta y) - f(x,y) = a\Delta{x} + b\Delta{y} + o(\rho)\]
    当$\Delta x \to 0,\Delta y \to 0$时，有：
    \[|f(x+\Delta x,y+\Delta y) - f(x,y)| \leqslant |a||\Delta x| + |b||\Delta y| + o(\rho) \to 0\]
    即表明其满足连续的条件：
    \[\lim_{\substack{\Delta x \to 0 \\ \Delta y \to 0}} f(x+\Delta x,y+\Delta y) = f(x,y)\]
    \circnum{1}证毕.注意到若取$\Delta y = 0$，有：
    \[f(x+\Delta x,y) - f(x,y) = a\Delta x + o(\Delta x)\]
    两边同除以$\Delta x$并取极限：
    \[\lim_{\Delta x \to 0}\frac{f(x+\Delta x,y)-f(x,y)}{\Delta x} = \lim_{\Delta x \to 0}\left(a+\frac{o(\Delta x)}{\Delta x}\right) = a = \pdv{f}{x}\]
    同理有$\displaystyle b =\pdv{f}{y}$，\circnum{2}证毕.
\end{proof}

与单变量情形不同的是，即使多变量函数$f(x,y)$的偏导数都存在，既不能保证$f$的连续性，也不能保证$f$的可微性.

现在研究加入一些条件限制后，偏导数与连续性和可微性的关系能否发生变化.

\begin{theorem}
    设函数$f(x,y)$在区域$D$内存在两个偏导数.
    \begin{enumerate}[circ]
        \item 如果偏导数$f'_x(x,y),f'_y(x,y)$在$D$内有界，那么$f$在$D$内连续.
        \item 如果偏导数$f'_x(x,y),f'_y(x,y)$在$D$内连续，那么$f$在$D$内可微.
    \end{enumerate}
\end{theorem}

\begin{proof}
    要证明$f(x,y)$连续，即要估计$|f(x+\Delta x,y+\Delta y) - f(x,y)|$的值：
    \[|f(x+\Delta x,y+\Delta y) - f(x,y)| \!\!\leqslant\!\! |f(x+\Delta x,y+\Delta y) - f(x+\Delta x,y)| + |f(x+\Delta x,y) - f(x,y)|\]
    上式中，利用Lagrange中值公式有：
    \[\begin{cases}
        |f(x+\Delta x,y+\Delta y) - f(x+\Delta x,y)| = |f'_y(x+\Delta x,y+\theta_1\Delta y)\Delta y|,\,0<\theta_1<1\\
        |f(x+\Delta x,y) - f(x,y)| = |f'_x(x+\theta_2\Delta x,y)\Delta x|,\,0<\theta_2<1
    \end{cases}\]
    若偏导数有界：$\begin{cases}
        |f'_x(x,y)|\leqslant M\\
        |f'_y(x,y)|\leqslant M
    \end{cases}$.则在$\Delta x \to 0,\Delta y \to 0$时有：
    \[|f(x+\Delta x,y+\Delta y) - f(x,y)| \leqslant M|\Delta y| + M|\Delta x| \to 0\]
    即表明在偏导数有界时连续，\circnum{1}得证.

    根据可微的定义，要证明可微性，即要证明在$\Delta x \to 0,\Delta y \to 0$时有：
    \[\lim_{\rho \to 0}\frac{|f(x+\Delta x,y+\Delta y) - f(x,y) - f'_x \Delta x - f'_y \Delta y|}{\rho} = 0\]
    注意到，根据Lagrange中值公式有：
    \begin{align*}
        &\frac{|f(x+\Delta x,y+\Delta y) - f(x,y) - f'_x \Delta x - f'_y \Delta y|}{\rho} \\
        \leqslant & \frac{|(f'_y(x+\Delta x,y+\theta_1\Delta y)-f'_y(x,y))\Delta y| + |(f'_x(x+\theta_2\Delta x,y)-f'_x(x,y))\Delta x|}{\rho} \\
        \leqslant & |f'_y(x+\Delta x,y+\theta_1\Delta y)-f'_y(x,y)| + |f'_x(x+\theta_2\Delta x,y)-f'_x(x,y)|
    \end{align*}
    若引入两个偏导数的连续性：
    \[\lim_{\substack{\Delta x \to 0 \\ \Delta y \to 0}}f'_y(x+\Delta x,y+\Delta y) = f'_y(x,y),\quad \lim_{\substack{\Delta x \to 0 \\ \Delta y \to 0}}f'_x(x+\Delta x,y+\Delta y) = f'_x(x,y)\]
    则有$|f'_y(x+\Delta x,y+\theta_1\Delta y)-f'_y(x,y)| + |f'_x(x+\theta_2\Delta x,y)-f'_x(x,y)| \to 0$成立，进而可微性得证，\circnum{2}证毕.
\end{proof}


\subsection{方向导数与梯度}
偏导数是沿着$x$或$y$方向在$f(x,y)$曲面上的导数，而对于其它方向上的导数，则引入方向导数的定义.

设$L$是过$(x,y)$的一条射线，其方向向量$\vb{e}$可以表示为：
\[\vb{e} = \cos\alpha\vb{i} + \cos\beta\vb{j},\quad\alpha+\beta=\frac{\pi}{2}\]
则有$L$上一动点$(x+t\cos\alpha,y+t\cos\beta)$，于是在$\vb{e}$方向上，$f(x,y)$变成了关于$t$的函数$f(t)$.
\begin{definition}
    如果下列极限存在，那么称极限值为$f$在$(x,y)$处沿方向$\vb{e}$的\overtext{方向导数}{directional derivative}，并记作：
    \[\pdv{f}{\vb{e}} \coloneq \lim_{t \to 0}\frac{f(x+t\cos\alpha,y+t\cos\beta)-f(x,y)}{t}\]
\end{definition}

\begin{theorem}
    设$f(x,y)$是平面区域$D$上的可微函数，则$f(x,y)$在$D$中任何一点沿任何方向$\vb{e}$的方向导数都存在，而且有：
    \[\pdv{f}{\vb{e}} = \pdv{f}{x}\cos\alpha + \pdv{f}{y}\cos\beta\]
\end{theorem}

\begin{proof}
    根据可微性有：
    \[f(x+\Delta x,y+\Delta y) - f(x,y) = f'_x(x,y)\Delta x + f'_y(x,y)\Delta y + o(\rho)\]
    特别取$\Delta x = t\cos\alpha,\Delta y = t\cos\beta$，则有$\rho = |t|$，上式两边同除以$t$取极限：
    \[\pdv{f}{\vb{e}} = f'_x\cos\alpha + f'_y\cos\beta\]
    由此得证.
\end{proof}

\begin{definition}
    称由函数的偏导数所决定的向量$\displaystyle\left(\pdv{f}{x},\pdv{f}{y}\right)$为$f(x,y)$在$(x,y)$处的\overtext{梯度}{gradient}，记作：
    \[\symbf{\nabla} f = \grad f \coloneq \left(\pdv{f}{x},\pdv{f}{y}\right) = \pdv{f}{x}\vb{i} + \pdv{f}{y}\vb{j}\]
\end{definition}

注意到，一个方向向量$\vb{e}$与一点梯度的内积为$\vb{e}$方向的方向导数：
\[\symbf{\nabla} f \cdot \vb{e} = \pdv{f}{\vb{e}} = |\symbf{\nabla}f|\cos\theta\]
也可以说，$f$在$(x,y)$上沿$\vb{e}$的方向导数是梯度在这个方向上的投影.这同时表明，沿着梯度的方向就是$f$在$(x,y)$附近变化率最大的方向.


\subsection{复合函数的微分和一阶微分形式的不变性}
\begin{theorem}
    设$z = f(\xi,\eta)$在$\tilde{D}$中可微；$\xi = \varphi(x,y),\eta = \psi(x,y)$在$D$中可微.则$z = f(\varphi(x,y),\psi(x,y))$作为$(x,y)\in\tilde{D}$中的函数也可微.
\end{theorem}

\begin{proof}
    根据$z = f(\xi,\eta)$的可微性，有：
    \[\Delta z = f(\xi+\Delta\xi,\eta+\Delta\eta) - f(\xi,\eta) = c_1\Delta\xi + c_2\Delta\eta + o(\tilde{\rho})\]
    其中$\tilde{\rho} = \sqrt{(\Delta\xi)^2 + (\Delta\eta)^2}$.又根据$\xi=\varphi(x,y),\eta=\psi(x,y)$的可微性：
    \begin{gather*}
        \Delta \xi = \varphi(x+\Delta x,y+\Delta y) - \varphi(x,y) = a_1\Delta x + a_2\Delta y + o(\rho)\\
        \Delta \eta = \psi(x+\Delta x,y+\Delta y) - \psi(x,y) = b_1\Delta x + b_2\Delta y + o(\rho)
    \end{gather*}
    其中$\rho = \sqrt{(\Delta x)^2 + (\Delta y)^2}$，将以上二式代入$\Delta z$的表达式中：
    \[\Delta z = (c_1 a_1 + c_2 b_1)\Delta x + (c_1 a_2 + c_2 b_2)\Delta y + c_1 o(\rho) + c_2 o(\rho) + o(\tilde{\rho})\]
    要证存在常数$A,B$使得$\Delta z =A\Delta x + B\Delta y + o(\rho)$，只需证明$\lim\limits_{\rho \to 0}\dfrac{o(\tilde{\rho})}{\rho} = 0$.
    根据$\tilde{\rho} = \sqrt{(\Delta\xi)^2 + (\Delta\eta)^2}$，有：
    \[\frac{\tilde{\rho}^2}{\rho^2} = \left(\frac{a_1\Delta x + a_2\Delta y + o(\rho)}{\rho}\right)^2 + \left(\frac{b_1\Delta x + b_2\Delta y + o(\rho)}{\rho}\right)^2\]
    注意到$\rho\to0$时有$\Delta x\to0,\Delta y\to0$，进而$\Delta\xi\to0,\Delta\eta\to0$，使得$\tilde{\rho}\to0$.且$|a_1|,|a_2|,|b_1|,|b_2|$小于$1$，故而$\frac{\tilde{\rho}}{\rho}$是有界的.从而有：
    \[\lim_{\rho\to0}\frac{o(\tilde{\rho})}{\rho} = \lim_{\rho\to0}\frac{o(\tilde{\rho})}{\tilde{\rho}}\cdot\frac{\tilde{\rho}}{\rho} = \lim_{\tilde{\rho}\to0}\frac{o(\tilde{\rho})}{\tilde{\rho}}\cdot\frac{\tilde{\rho}}{\rho} = 0\]
    由此得证.
\end{proof}
在上述推导中，不难得到以下两个结论：
\begin{theorem}
    \textbf{一阶微分形式的不变性}

    对于函数$z = f(\xi,\eta) = g(x,y)$，$z$的微分关于两组自变量$(\xi,\eta)$和$(x,y)$的微分形式是相等的，即：
    \[\dd{z} = \pdv{z}{\xi}\dd{\xi} + \pdv{z}{\eta}\dd{\eta} = \pdv{z}{x}\dd{x} + \pdv{z}{y}\dd{y}\]
\end{theorem}
\begin{theorem}
    \textbf{链式法则}

    对复合函数$f(\xi,\eta),\xi = \varphi(x,y),\eta = \psi(x,y)$，有：
    \[\pdv{f}{x} = \pdv{f}{\xi}\pdv{\xi}{x} + \pdv{f}{\eta}\pdv{\eta}{x},\qquad \pdv{f}{y} = \pdv{f}{\xi}\pdv{\xi}{y} + \pdv{f}{\eta}\pdv{\eta}{y}\]
\end{theorem}


\subsection{向量值函数的微商和微分}
\begin{definition}
    若下列极限存在，则称之为向量值函数$\vb{r}(t)$关于$t$的导数：
    \[\vb{r}'(t) = \dv{\vb{r}(t)}{t} \coloneq \lim_{\Delta t \to 0}\frac{\vb{r}(t+\Delta t) - \vb{r}(t)}{\Delta t}\]
\end{definition}
从几何上看，$\vb{r}(t)$在$t_0$的导数$\vb{r}'(t)$就是参数曲线在对应$t_0$一点上的切向量，指向$t$增加的方向.

引进向量的微分：
\[\dd{\vb{r}(t)} = \dv{\vb{r(t)}}{t}\dd{t} \quad\text{或}\quad \dd{\vb{r}} = \dd{x}\vb{i} + \dd{y}\vb{j} + \dd{z}\vb{k}\]
向量值函数微分的向量运算满足如下性质：
\begin{enumerate}[property]
    \item $\displaystyle\dv{t}(f\vb{a}) = f\dv{\vb{a}}{t} + \dv{f}{t}\vb{a}$
    \item $\displaystyle\dv{t}(\vb{a \cdot b}) = \dv{\vb{a}}{t}\cdot\vb{b} + \vb{a}\cdot\dv{\vb{b}}{t}$
    \item $\displaystyle\dv{t}(\vb{a \times b}) = \dv{\vb{a}}{t}\times\vb{b} + \vb{a}\times\dv{\vb{b}}{t}$
\end{enumerate}

接下来在向量值函数的基础上讨论求导与微分在更一般函数上的情形.

设有映射$\vb{f}:\mathbb{R}^n \to \mathbb{R}^m$，即$n$维向量$\vb{x}$到$m$维向量$\vb{y}$的映射：
\[\vb{y} = (y_1,y_2,\cdots,y_m) = (f_1(\vb x),f_2(\vb x),\cdots,f_m(\vb x))\]
定义对$\vb{f}$的微分就是对每一个分量的微分：
\[\dd{\vb{f(x)}} = (\dd{f_1(\vb{x})},\dd{f_2(\vb{x})},\cdots,\dd{f_m(\vb{x})}),\quad\dd{f_j} = \sum_{i=1}^{n}\pdv{f_j}{x_i}\dd{x_i},\quad j = 1,2,\cdots,m\]
上述运算写为矩阵形式就是：
\[\begin{bmatrix}
    \dd{f_1} \\ \vdots \\ \dd{f_m}
\end{bmatrix} = \begin{bmatrix}
    \displaystyle\pdv{f_1}{x_1} & \cdots & \displaystyle\pdv{f_1}{x_n} \\
    \vdots & & \vdots \\
    \displaystyle\pdv{f_m}{x_1} & \cdots & \displaystyle\pdv{f_m}{x_n} \\
\end{bmatrix} \begin{bmatrix}
        \dd{x_1} \\ \vdots \\ \dd{x_n}
\end{bmatrix}\]
为了进一步简化计算，作如下定义：
\begin{definition}
    称下列矩阵为向量值函数$\vb{f(\vb{x})}$的\textbf{Jacobi矩阵}：
    \[\vb{J}_{\vb x}(\vb f) = \begin{bmatrix}
    \displaystyle\pdv{f_1}{x_1} & \cdots & \displaystyle\pdv{f_1}{x_n} \\
    \vdots & & \vdots \\
    \displaystyle\pdv{f_m}{x_1} & \cdots & \displaystyle\pdv{f_m}{x_n} \\
    \end{bmatrix}_{m \times n}\]
    特别地，当$m=n$时，上述矩阵为方阵，其行列式称为\textbf{Jacobi行列式}，记作：
    \[\det\vb{J}_{\vb{x}}(\vb{f}) = \pdv{(y_1,y_2,\cdots,y_n)}{(x_1,x_2,\cdots,x_n)}\]
\end{definition}
从几何上看，Jacobi行列式就是经过同维度映射$\vb{f}$后微元的缩放因子.

Jacobi行列式是同维度向量值函数复合求导的基本计算手段.记有两个映射：
\[\begin{cases}
    (x,y)\mapsto(u,v):& u = u(x,y),\quad v = v(x,y)\\
    (u,v)\mapsto(\xi,\eta):& \xi = \xi(u(x,y),v(x,y)),\quad \eta = \eta(u(x,y),v(x,y))
\end{cases}\]
它们复合后的映射为$(x,y)\mapsto(\xi,\eta)$，$\xi$和$\eta$分别对$x$和$y$的偏导数就是复合函数的偏导数，这个过程中的Jacobi行列式满足：
\begin{theorem}
    \textbf{Jacobi行列式的链式法则}
    \[\pdv{(\xi,\eta)}{(x,y)} = \pdv{(\xi,\eta)}{(u,v)}\pdv{(u,v)}{(x,y)}\]
\end{theorem}


\section{隐函数定理和逆映射定理}
若$\mathbb{R}^2$中的点$(x,y)$满足$F(x,y) = 0$，则$F$给出了$\mathbb{R}^2$上的一条曲线，或一种函数关系.

例如，$F(x,y) = ax + by + c = 0$给出了一条直线.这条直线可以显式地表示为：
\[y = -\frac{1}{b}(ax + c),\quad b \neq 0 \quad\text{或}\quad x = -\frac{1}{a}(by + c),\quad a \neq 0\]

而$F(x,y) = x^2 + y^2 - 1 = 0$给出了一个单位圆.要想将其显式地表示，需要将其分为上下半圆两个方程：
$\begin{cases}
    y = \sqrt{1-x^2} \\
    y = -\sqrt{1-x^2}
\end{cases}$，特别注意到它们的导数$y' = \dfrac{\pm x}{\sqrt{1-x^2}}$中$x$不能取到$\pm 1$，则圆左右两端点的切线无法被描述.为了完整的描述这个圆的性质，在前述两个方程的基础上还要再引入两个方程$\begin{cases}
    x = \sqrt{1-y^2} \\
    x = -\sqrt{1-y^2}
\end{cases}$，总共四个方程才能完整显式地描述.

又例如方程$F(x,y) = \sin(x+y) + 2x + y = 0$根本无法得出显表示；$F(x,y) = x^2 + y^2 + 1 = 0$根本没有实数解.

因此，不能总是期待得到显表示，或是从显表示中得到函数的可微性，这就是研究隐函数的必要性.

\subsection{隐函数的存在性和微商}
假如方程$F(x,y) = 0$在$(x_0,y_0)$邻域内存在隐函数，不妨设为$y = f(x)$，则有$F(x,f(x)) \equiv 0$，两边求导，得到隐函数的导数：
\[F'_x + F'_y f'(x) = 0 \quad\Rightarrow\quad f'(x) = \dv{y}{x} = -\frac{F'_x(x,y)}{F'_y(x,y)}\]
上述分析表明，即使不能得到隐函数的显表示，求取以$x$和$f(x)$表示的导数仍然是可行的.

\begin{theorem}
    \textbf{隐函数存在定理}
    \label{thm:ExistenceOfImplicitFunctions}

    设开区域$D \subset \mathbb{R}^2$，$(x_0,y_0) \in D$.如果$F(x,y)$在$D$中有定义且满足：
    \begin{enumerate}[circ]
        \item $F(x,y) \in C^1(D)$，即$F$在$D$中有连续的一阶偏导数.
        \item $F(x_0,y_0) = 0$.
        \item $F'_y(x_0,y_0) \neq 0$.
    \end{enumerate}
    那么存在$(x_0,y_0)$的一个矩形开邻域$(a,b) \times (c,d) \subset D$，使得对每一个$x \in (a,b)$，方程$F(x,y) = 0$存在唯一解$y = f(x)$，满足$y_0 = f(x_0)$且$y = f(x)$连续并有下列微商：
    \[f'(x) = \dv{y}{x} = -\frac{F'_x(x,y)}{F'_y(x,y)}\]
    特别地，若条件\circnum{3}改为$F'_x(x_0,y_0) \neq 0$，则存在隐函数$x = g(y)$具有上述类似特性.
\end{theorem}

曲线$F(x,y) = 0$在$(x_0,y_0)$的邻域内有隐函数$y = f(x)$的几何含义是：在$(x_0,y_0)$邻域内曲线$F(x,y) = 0$可以用$y = f(x)$表示.

\begin{figure}[H]
    \centering
    \vspace{-4em}
    \includegraphics[width=\textwidth]{existence-of-inplicit-func/fig.pdf}
    \vspace{-4em}
    \caption{隐函数的几何含义}
\end{figure}

\begin{proof}
    \textbf{存在性：}
    
    根据条件\circnum{3}有$F'_y(x_0,y_0) \neq 0$，不妨设$F'_y(x_0,y_0) > 0$.则根据连续的保号性，存在$(x_0,y_0)$的邻域$[a',b'] \times [c,d] \subset D$，使得$F'_y(x,y)$在其上大于$0$.

    换言之，固定一个$\forall x \in [a',b']$，$F(x,y)$作为关于$y$的函数由于其导数严格大于$0$，$F(x,y)$在$y \in [c,d]$上严格单调增.

    特别地，取$x_0 \in [a',b']$，有$F(x_0,y)$在$y \in [c,d]$上严格单调增，则有：
    \[F(x_0,c) < F(x_0,y_0) < F(x_0,d)\]
    根据条件\circnum{2}，有$F(x_0,y_0) = 0$，则有：
    \[F(x_0,c) < 0,\qquad F(x_0,d) > 0\]
    由$x$在$[a',b']$上的任意性，上述结论可以推广到$x_0$的邻域$(a,b)$上，即：
    \[F(x,c) < 0 < F(x,d),\quad x \in (a,b) \subset [a',b']\]
    根据介值定理和$F(x_0,y)$的严格单调增，对于一个固定的$x \in (a,b)$，在$(c,d)$中存在唯一的$y = f(x)$使得$F(x,f(x)) = 0$.即$y = f(x)$是满足$F(x,y) = 0$的隐函数，且有$F(x_0,y_0) = 0$.

    \textbf{可微性：}

    对$\forall x \in (a,b)$，取$\Delta x$使得$(x + \Delta x) \in (a,b)$，并且记$\Delta y = f(x+\Delta x) - f(x)$.

    $(x,f(x))$和$(x+\Delta x,f(x+\Delta x))$仍在隐函数存在的邻域内，因此有：
    \[F(x,f(x)) = F(x+\Delta x,f(x+\Delta x)) = F(x+\Delta x,y+\Delta y) = 0\]
    于是有：
    \begin{align*}
        0 =& F(x+\Delta x,y+\Delta y) - F(x,y) \\
        =& (F(x+\Delta x,y+\Delta y) - F(x+\Delta x,y)) + (F(x+\Delta x,y) - F(x,y))
    \end{align*}
    根据Lagrange中值公式，有$0<\theta_1<1,0<\theta_2<1$使得：
    \begin{gather*}
        0 = F'_y(x+\Delta x,y+\theta_1\Delta y)\Delta y + F'_x(x+\theta_2\Delta x,y)\Delta x \\
        \Delta y = -\frac{F'_x(x+\theta_2\Delta x,y)}{F'_y(x+\Delta x,y+\theta_1\Delta y)}\Delta x
    \end{gather*}
    根据条件\circnum{1}，$F'_x,F'_y$在$D$中连续，同时根据条件\circnum{3}分母不为$0$，因而存在一个$(x_0,y_0)$的邻域使得有界性成立：
    \[|\Delta y| = \frac{|F'_x(x+\theta_2\Delta x,y)|}{|F'_y(x+\Delta x,y+\theta_1\Delta y)|}|\Delta x| \leqslant M|\Delta x|\]
    上式表明，当$\Delta x \to 0$时，有$|\Delta y| = |f(x+\Delta x) - f(x)| \to 0$，即$y = f(x)$连续.于是有：
    \[\dv{y}{x} = \lim_{\Delta x \to 0}\frac{\Delta y}{\Delta x} = \lim_{\Delta x \to 0}\left(-\frac{F'_x(x+\theta_2\Delta x,y)}{F'_y(x+\Delta x,y+\theta_1\Delta y)}\right) = -\frac{F'_x(x,y)}{F'_y(x,y)}\]
\end{proof}

\begin{mistake}
    关于\cref{thm:ExistenceOfImplicitFunctions}\textbf{(隐函数存在定理)}的条件\circnum{3}：$F'_y(x_0,y_0)\neq 0$只是隐函数存在的充分条件，而不是必要条件.

    也就是说，就算$F'_y(x_0,y_0) = 0$，也可能有隐函数.
\end{mistake}

\begin{theorem}
    若一般曲线$F(x,y) = 0$在$(x_0,y_0)$处自相交，则其在这一点的两个偏导数必定为$0$.
\end{theorem}

\begin{proof}
    自相交则表明$(x_0,y_0)$的邻域内不存在隐函数，因为不满足函数的单值条件.不存在隐函数，则\cref{thm:ExistenceOfImplicitFunctions}\textbf{(隐函数存在定理)}中的三个条件至少有一个为假，显然对于一般曲线\circnum{1}和\circnum{2}都为真，故而：
    \[F'_x(x_0,y_0) = F'_y(x_0,y_0) = 0\]
\end{proof}

前述推导都是二维空间中的一个二维方程，接下来对隐函数的存在条件和隐函数的计算方法做一定的推广.
\subsubsection{三维空间中的一个方程}
三维空间中有一个方程意味着三个变量具有一个约束，即其中一个变量由另外两个变量决定，下列方程隐函数是二元的，且隐函数只有一个.
\[F(x,y,z) = 0\]
如果存在$(x_0,y_0,z_0)$使得$F(x_0,y_0,z_0) = 0$，同时在$(x_0,y_0,z_0)$邻域内存在隐函数$z = f(x,y)$，则有：
\[F(x,y,f(x,y)) \equiv 0\]
对上式分别对$x,y$求导，由此解得隐函数$z = f(x,y)$的两个偏导数：
\[\pdv{z}{x} = -\frac{F'_x(x,y,z)}{F'_z(x,y,z)},\quad\pdv{z}{y} = -\frac{F'_y(x,y,z)}{F'_z(x,y,z)}\]
上述偏导数成立的条件是$F'_z(x_0,y_0,z_0) \neq 0$.同理，若取隐函数为$x = x(y,z)$和$y = y(x,z)$，则也有相应的要求.由于隐函数的取法具有任意性，所以要想使隐函数存在，要满足在$(x_0,y_0,z_0)$上$F'_x,F'_y,F'_z$都不为$0$，即：
\[\grad F(x_0,y_0,z_0) \neq \bm{0}\]
总而言之，隐函数存在的一个条件是有连续偏导数、有一个解且该点处梯度不为$\bm{0}$.


\subsubsection{三维空间中两个方程构成的方程组}
对以下方程，意味着有三个变量两个约束，隐函数是单变量的，且隐函数有两个.
\[\begin{cases}
    F(x,y,z) = 0 \\
    G(x,y,z) = 0
\end{cases}\]
如果上述方程组有一个解$(x_0,y_0,z_0)$且在这个解的邻域内有隐函数：
\[\begin{cases}
    y = y(x)\\
    z = z(x)
\end{cases} \Rightarrow\quad \begin{cases}
    F(x,y(x),z(x)) = 0\\
    G(x,y(x),z(x)) = 0\\
\end{cases}\]
对$x$求导，得到两个隐函数关于$x$的导数：
\[y'(x) = -\frac{F'_x G'_z - F'_z G'_x}{F'_y G'_z - F'_z G'_y},\qquad z'(x) = \frac{F'_x G'_y - F'_y G'_x}{F'_y G'_z - F'_z G'_y}\]
上述导数存在的条件是关于$y,z$的Jacobi行列式不为零：
\[\pdv{(F,G)}{(y,z)} = \begin{vmatrix}
    \displaystyle\pdv{F}{y} & \displaystyle\pdv{F}{z} \\
    \displaystyle\pdv{G}{y} & \displaystyle\pdv{G}{z} \\
\end{vmatrix} \neq 0\]
若取其它隐函数也有类似的要求：
\[\pdv{(F,G)}{(x,y)} \neq 0,\quad \pdv{(F,G)}{(y,z)} \neq 0,\quad\pdv{(F,G)}{(x,z)} \neq 0 \quad\Leftrightarrow\quad \grad F \times \grad G \neq 0\]
于是隐函数存在的一个条件是均有连续偏导数，有一个解，在这个解处两个梯度的叉乘不为零.


\subsubsection{四维空间中两个方程构成的方程组}
对方程：
\[\begin{cases}
    F(x,y,u,v) = 0 \\
    G(x,y,u,v) = 0
\end{cases}\]
有四个变量两个约束，有两个二元隐函数.其它推导过程与条件同上.


\subsection{逆映射的微商}
在一元函数中，反函数存在的条件是$y=f(x)$有连续的导函数，且$f'(x_0)\neq 0$.它此前的证明是通过$f'(x)$连续所具有的局部保号性，推出$x_0$邻域内$f(x)$的严格单调性，进而证明反函数$x = f^{-1}(y)$存在且等于$[f'(x)]^{-1}$.

利用隐函数存在定理，可以从另一个角度给出证明.

\begin{proof}
    已知$f'(x_0)\neq 0$，令$F(x,y) = y - f(x)$，考虑方程：
    \[F(x,y) = y - f(x) = 0\]
    注意到有$F'_x(x_0,y_0) = f'(x_0) \neq 0$，根据隐函数存在定理，存在隐函数$x = x(y) = f^{-1}(y)$，同时得到其导数为：
    \[\dv{x}{y} = \dv{f^{-1}}{y} = -\frac{F'_y}{F'_x} = \frac{1}{f'(x)} \quad\Rightarrow\quad \dv{f^{-1}}{y}\dv{f}{x} = 1\]
\end{proof}

\begin{theorem}
    \textbf{逆映射定理}

    设映射$(x,y)\mapsto(u,v):\begin{cases}
        u = u(x,y) \in C^1 \\
        v = v(x,y) \in C^1
    \end{cases}$，若在一点$(x_0,y_0)$处有$\displaystyle\pdv{(u,v)}{(x,y)}\neq 0$，则在$(u_0,v_0) = (u(x_0,y_0),v(x_0,y_0))$的邻域内存在逆映射$(u,v)\mapsto(x,y):\begin{cases}
        x = x(u,v)\\
        y = y(u,v)
    \end{cases}$，且$x(u,v),y(u,v)$均可导，其偏导数的Jacobi行列式满足：
    \[\pdv{(x,y)}{(u,v)}\cdot\pdv{(u,v)}{(x,y)} = 1\]
\end{theorem}

逆映射定理所要求的Jacobi行列式不为零，意味着正映射过程中不得发生降维，否则造成了信息丢失，逆映射无法确立.


\subsection{从微分看隐函数定理}
根据\cref{thm:ExistenceOfImplicitFunctions}\textbf{(隐函数存在定理)}中的条件\circnum{1}：$F(x,y)\in C^1(D)$，因而有：
\[\dd{F} = F'_x\dd{x} + F'_y\dd{y}\]
加入条件\circnum{2}：$F(x,y)=0$，若能给出隐函数$y = f(x)$，则在隐函数生效的邻域内有：
\[F(x,f(x))\equiv 0 \quad\Rightarrow\quad \dd{F} = F'_x\dd{x} + F'_y\dd{y} = 0\] 
同理若给出的隐函数为$x = g(y)$，同样有上式.对于上式，只要有一个系数不为零，即满足隐函数存在定理的条件\circnum{3}，不妨设$F'_y \neq 0$，则可以得到隐函数的微分：
\[\dd{y} = f'(x)\dd{x} = -\frac{F'_x(x,f(x))}{F'_y(x,f(x))}\dd{x}\]
从上述推导推广至$n$个变量、一个方程约束的情形：
\begin{gather*}
    F(x_1,x_2,\cdots,x_n) = 0,\quad (x_1,x_2,\cdots,x_n) \in D \subset \mathbb{R}^n\\
    \dd{F} = F'_{x_1}\dd{x_1} + F'_{x_2}\dd{x_2} + \cdots + F'_{x_n}\dd{x_n} = 0
\end{gather*}
其中，不妨设$F'_{x_n} \neq 0$，则有：
\[\dd{x_n} = -\frac{1}{F'_{x_n}}(F'_{x_1}\dd{x_1} + \cdots + F'_{x_{n-1}}\dd{x_{n-1}}) \quad\Rightarrow\quad \pdv{x_n}{x_j} = -\frac{F'_{x_j}}{F'_{x_n}},\quad j = 1,\cdots,n-1\]
在上述推导中可以发现两个事实：
\begin{itemize}
    \item 约束方程的微分可以写为各自变量微分的线性组合.
    \item 某个自变量的微分可以写为剩余自变量微分的线性组合，且各项系数是与之对应的偏导数.
\end{itemize}
接下来，进一步推广至$n$个变量，$m$个约束方程：
\[F_i(x_1,\cdots,x_n) \in C^1(D),\quad i = 1,2,\cdots,m \quad\Rightarrow\quad \begin{cases}
    F_1(x_1,\cdots,x_n) = 0\\
    F_2(x_1,\cdots,x_n) = 0\\
    \vdots\\
    F_m(x_1,\cdots,x_n) = 0
\end{cases}\]
对上述方程组微分，得到如下线性方程组：
\[\begin{cases}
    \displaystyle\dd{F_1} = \pdv{F_1}{x_1}\dd{x_1} + \cdots + \pdv{F_1}{x_n}\dd{x_n} = 0 \\[1em]
    \displaystyle\dd{F_2} = \pdv{F_2}{x_1}\dd{x_1} + \cdots + \pdv{F_2}{x_n}\dd{x_n} = 0 \\[1em]
    \vdots\\
    \displaystyle\dd{F_m} = \pdv{F_m}{x_1}\dd{x_1} + \cdots + \pdv{F_m}{x_n}\dd{x_n} = 0
\end{cases}\]
将上述线性方程组移项并写成矩阵形式，不妨设$m < n$：
\begin{align*}
    \begin{bmatrix}
        \displaystyle\pdv{F_1}{x_1} & \cdots & \displaystyle\pdv{F_1}{x_m} \\
        \vdots & & \vdots\\
        \displaystyle\pdv{F_m}{x_1} & \cdots & \displaystyle\pdv{F_m}{x_m} \\
    \end{bmatrix}\begin{bmatrix}
        \dd{x_1} \\ \vdots \\ \dd{x_m}
    \end{bmatrix} &= \begin{bmatrix}
        \displaystyle -\pdv{F_1}{x_{m+1}}\dd{x_{m+1}} -\cdots -\pdv{F_1}{x_{n}}\dd{x_{n}}\\
        \vdots \\
        \displaystyle -\pdv{F_m}{x_{m+1}}\dd{x_{m+1}} -\cdots -\pdv{F_m}{x_{n}}\dd{x_{n}}\\
    \end{bmatrix}\\
    \underbrace{\begin{bmatrix}
        \displaystyle\pdv{F_1}{x_1} & \cdots & \displaystyle\pdv{F_1}{x_m} \\
        \vdots & & \vdots\\
        \displaystyle\pdv{F_m}{x_1} & \cdots & \displaystyle\pdv{F_m}{x_m} \\
    \end{bmatrix}}_{\vb{J}_1} \begin{bmatrix}
        \dd{x_1} \\ \vdots \\ \dd{x_m}
    \end{bmatrix} &= -\underbrace{\begin{bmatrix}
        \displaystyle\pdv{F_1}{x_{m+1}} & \cdots & \displaystyle\pdv{F_1}{x_n} \\
        \vdots & & \vdots\\
        \displaystyle\pdv{F_m}{x_{m+1}} & \cdots & \displaystyle\pdv{F_m}{x_n} \\
    \end{bmatrix}}_{\vb{J}_2}\begin{bmatrix}
        \dd{x_{m+1}} \\ \vdots \\ \dd{x_n}
    \end{bmatrix}
\end{align*}
记上述方程中左边的Jacobi矩阵为$\vb{J}_1$，右边的为$\vb{J}_2$.
求解上述方程，首先应当保证$\det \vb{J}_1 \neq 0$，这样才能求逆.
\begin{gather*}
    \begin{bmatrix}
        \dd{x_1} \\ \vdots \\ \dd{x_m}
    \end{bmatrix} = -\vb{J}_1^{-1} \vb{J}_2 \begin{bmatrix}
        \dd{x_{m+1}} \\ \vdots \\ \dd{x_n}
    \end{bmatrix} = \begin{bmatrix}
    \displaystyle\pdv{x_1}{x_{m+1}} & \cdots & \displaystyle\pdv{x_1}{x_n} \\
    \vdots & & \vdots \\
    \displaystyle\pdv{x_m}{x_{m+1}} & \cdots & \displaystyle\pdv{x_m}{x_n} \\
\end{bmatrix} \begin{bmatrix}
        \dd{x_{m+1}} \\ \vdots \\ \dd{x_n}
    \end{bmatrix}
\end{gather*}
显然，上式中的$x_1,\cdots,x_m$都是关于$x_{m+1},\cdots,x_n$的隐函数，记这些隐函数的偏导数构成的Jacobi矩阵为$\mathcal{D}\vb{f}$
\[\begin{cases}
    x_1 = f_1(x_{m+1},\cdots,x_n) \\
    \vdots \\
    x_m = f_m(x_{m+1},\cdots,x_n) \\
\end{cases} \Rightarrow\quad \mathcal{D}\vb{f} \coloneq \begin{bmatrix}
    \displaystyle\pdv{f_1}{x_{m+1}} & \cdots & \displaystyle\pdv{f_1}{x_n} \\
    \vdots & & \vdots \\
    \displaystyle\pdv{f_m}{x_{m+1}} & \cdots & \displaystyle\pdv{f_m}{x_n} \\
\end{bmatrix} = \begin{bmatrix}
    \displaystyle\pdv{x_1}{x_{m+1}} & \cdots & \displaystyle\pdv{x_1}{x_n} \\
    \vdots & & \vdots \\
    \displaystyle\pdv{x_m}{x_{m+1}} & \cdots & \displaystyle\pdv{x_m}{x_n} \\
\end{bmatrix}\]
结合上述二式，得到如下结论：
\[\mathcal{D}\vb{f} = -\vb{J}_1^{-1}\vb{J}_2,\quad \det \vb{J}_1 \neq 0\]
其中$\vb{J}_1$是$F$对设为隐函数的变量的Jacobi矩阵，$\vb{J}_2$是$F$对其余变量的Jacobi矩阵.
这就是求解多变量函数隐函数偏导数的方法.二元函数的结论$\dv{y}{x} = -\frac{F'_x}{F'_y}$就是上述公式的退化形式.


\section{空间曲线与曲面}
用点的集合能够表示空间中的一个几何对象.例如：
\[\begin{cases}
    \text{平面曲线：}\{(x,y)\,|\,y = f(x)\} \Leftrightarrow \vb{r}(x) = (x,f(x)) = x\vb{i} + f(x)\vb{j} \\
    \text{空间曲面：}\{(x,y,z)\,|\,z = f(x,y)\} \Leftrightarrow \vb{r}(x,y) = (x,y,f(x,y)) = x\vb{i} + y\vb{j} + f(x,y)\vb{k} \\
\end{cases}\]
转为参数表示则为：
\[\begin{cases}
    \text{参数曲线：}\{(x(t),y(t),z(t))\,|\,t\in[\alpha,\beta]\}\subset\mathbb{R}^3 \\
    \text{参数曲面：}\{(x(u,v),y(u,v),z(u,v))\,|\,(u,v)\in D \subset\mathbb{R}^2\}\subset\mathbb{R}^3
\end{cases}\]
隐表示的平面曲线则为$F(x,y)=0$，空间曲面$F(x,y,z)=0$，空间曲线则以两个空间曲面联立消自由度，表示为$\begin{cases}
    F(x,y,z) = 0\\
    G(x,y,z) = 0
\end{cases}$.


\subsection{参数曲线}
\subsubsection{切向量与切线方程}
已知参数曲线在一点$M(x(t_0),y(t_0),z(t_0))$的切向量为：
\[\vb{r}'(t_0) = x'(t_0)\vb{i} + y'(t_0)\vb{j} + z'(t_0)\vb{k}\]
对应的点向式切线方程为：
\[\frac{x - x(t_0)}{x'(t_0)} = \frac{y - y(t_0)}{y'(t_0)} = \frac{z - z(t_0)}{z'(t_0)}\]
对应的法平面方程为：
\[x'(t_0)(x-x(t_0)) + y'(t_0)(y-y(t_0)) + z'(t_0)(z-z(t_0)) = 0\]
\begin{definition}
    对于参数曲线，若有$x'(t),y'(t),z'(t)$连续，那么称曲线是\overtext{光滑曲线}{smooth curve}.当曲线连续且分成有限光滑曲线段，则称曲线是\overtext{逐段光滑}{piecewise smooth}的.当曲线上对于$t_0$一点有$|\vb{r}'(t_0)|\neq0$，则称该点为曲线的\overtext{正则点}{regular point}，否则称为\overtext{奇点}{singular point}.如果曲线上的所有点都是正则点，则称该曲线为\overtext{正则曲线}{regular curve}.
\end{definition}


\subsubsection{弧长}
对光滑参数曲线$\vb{r}(t)$的参变量$t$做分割：
\[T:\alpha = t_0 < t_1 < \cdots < t_n = \beta\]
则曲线的总长度为：
\begin{align*}
    l(T) &= \sum_{i=1}^{n}|\vb{r}(t_i) - \vb{r}(t_{i-1})| \\
    &= \sum_{i=1}^{n}\sqrt{[x(t_i)-x(t_{i-1})]^2 + [y(t_i)-y(t_{i-1})]^2 + [z(t_i)-z(t_{i-1})]^2}
\end{align*}
根据Lagrange微分中值公式，取$t_{i-1} < \xi_i,\eta_i,\zeta_i < t_i$，$\Delta t_i = t_i - t_{i-1}$：
\[l(T) = \sum_{i=1}^{n}\sqrt{{x'}^2(\xi_i) + {y'}^2(\eta_i) + {z'}^2(\zeta_i)}\Delta t_i\]
注意到，上式不是标准的Riemann和形式，应当为其构造一个近似的Riemann和.

对于光滑曲线，$x'(t),y'(t),z'(t)$是连续的，故而：
\[f(\xi,\eta,\zeta) = \sqrt{{x'}^2(\xi) + {y'}^2(\eta) + {z'}^2(\zeta)}\]
是定义在三维闭区间$[\alpha,\beta]^3$上的三元连续函数，进而其一致连续，即$\forall\varepsilon>0$，$\exists\delta_1>0$，使得在$\left\|T\right\|<\delta_1$时，对任意$(\xi_i,\eta_i,\zeta_i)\in[t_{i-1},t_i]^3$都有：
\begin{align*}
    \left|\sqrt{{x'}^2(\xi_i) + {y'}^2(\eta_i) + {z'}^2(\zeta_i)} - \sqrt{{x'}^2(t_i) + {y'}^2(t_i) + {z'}^2(t_i)}\right| &< \varepsilon \\
    \left|f(\xi_i,\eta_i,\zeta_i) - f(t_i)\right| &< \varepsilon \\
    \left|\sum_{i=1}^{n}f(\xi_i,\eta_i,\zeta_i)\Delta t_i - \sum_{i=1}^{n}f(t_i)\Delta t_i\right| &< \varepsilon\sum_{i=1}^{n}\Delta t_i
\end{align*}
上式中，$f(t_i)$因$x'(t),y'(t),z'(t)$的连续而连续，根据\cref{thm:ClassOfIntegrableFunctions}\textbf{(连续必可积)}和Riemann可积的定义：对$\forall\varepsilon>0$，$\exists\delta_2>0$使得在$\left\|T\right\|<\delta_2$时有：
\[\left|\sum_{i=1}^{n}f(t_i)\Delta t_i - \int_{\alpha}^{\beta}f(t)\dd{t}\right| < \varepsilon\]
于是当$\left\|T\right\|<\delta = \min\{\delta_1,\delta_2\}$时，有：
\begin{align*}
    \left|\sum_{i=1}^{n}f(\xi_i,\eta_i,\zeta_i)\Delta t_i - \sum_{i=1}^{n}f(t_i)\Delta t_i\right| + \left|\sum_{i=1}^{n}f(t_i)\Delta t_i - \int_{\alpha}^{\beta}f(t)\dd{t}\right| &< \varepsilon\sum_{i=1}^{n}\Delta t_i + \varepsilon
\end{align*}
对于左边又有：
\[\left|l(T) - \int_{\alpha}^{\beta}\!\!f(t)\dd{t}\right| \leqslant \left|\sum_{i=1}^{n}f(\xi_i,\eta_i,\zeta_i)\Delta t_i - \sum_{i=1}^{n}f(t_i)\Delta t_i\right| + \left|\sum_{i=1}^{n}f(t_i)\Delta t_i - \int_{\alpha}^{\beta}\!\!f(t)\dd{t}\right|\]
因而将光滑曲线$L$的弧长$l$定义为：
\[l \coloneq \lim_{\left\|T\right\|\to0}l(T) = \int_{\alpha}^{\beta}f(t)\dd{t} = \int_{\alpha}^{\beta}|\vb{r}'(t)|\dd{t} = \int_{\alpha}^{\beta}\sqrt{{x'}^2(t) + {y'}^2(t) + {z'}^2(t)}\dd{t}\]
特别地，对于逐段光滑曲线，其弧长可以上式作分段积分求和.若是平面曲线，取$z(t)\equiv0$即可.若平面曲线由显表示$y=f(x)$给出，上式可简化为：
\[l \coloneq \int_{a}^{b}\sqrt{1+{f'}^2(x)}\dd{x},\quad x\in[a,b]\]


\subsubsection{弧长参数}
设正则曲线$L$上有一从起点$A$出发的动点$M(t)$，则$A$到$M(t)$弧长为：
\[s(t) = \int_{\alpha}^{t}|\vb{r}'(\tau)|\dd{\tau},\quad t\in[\alpha,\beta]\]
上式是一个变上限积分给出的函数$s(t)$，根据正则曲线的性质，$s(t)$是严格单调增的，因而存在反函数$t(s)$.于是，曲线也可以以弧长为参数表达：
\[\vb{r} = \vb{r}(t) = \vb{r}(s),\quad x=x(s),y=y(s),z=z(s)\]
由于弧长作为曲线自身固有的几何量，不受外部坐标变化的影响，因此将以$s$为参数的参数曲线方程称为\textbf{自然方程}.

根据一阶微分形式的不变性，有：
\[\dd{\vb{r}} = \vb{r}'(t)\dd{t} = \vb{r}'(s)\dd{s}\]
根据反函数的性质，有：
\[\dv{s}{t} = |\vb{r}'(t)|,\quad \dv{t}{s} = \frac{1}{|\vb{r}'(s)|} \quad\Rightarrow\quad \vb{r}'(s) = \vb{r}'(t)\dv{s}{t} = \frac{\vb{r}'(t)}{|\vb{r}'(t)|}\]
即以弧长为参数的切向量是单位向量，进而推出：
\[|\vb{r}'(s)|^2 = \left(\dv{x}{s}\right)^2 + \left(\dv{y}{s}\right)^2 + \left(\dv{z}{s}\right)^2 = 1 \quad\Rightarrow\quad (\dd{s})^2 = (\dd{x})^2 + (\dd{y})^2 + (\dd{z})^2\]
上述结论可以看作是微元形式的勾股定理.


\subsubsection{曲率}
在物理学中，习惯上认为$\dot{x}$是$x$对时间$t$的导数.但是在数学的习惯上，记：
\[\dot{\vb{r}} = \dv{\vb{r}}{s},\quad \ddot{\vb{r}} = \dv[2]{\vb{r}}{s},\quad \vb{r}' = \dv{\vb{r}}{t},\quad \vb{r}'' = \dv[2]{\vb{r}}{t}\]
\begin{definition}
    假定对正则曲线$L$有连续的二阶导数.由于切向量$\dot{\vb{r}}$是单位向量，总是有$\dot{\vb{r}}\cdot\dot{\vb{r}} = 1$，从而有：
    \[\dv{s}(\dot{\vb{r}}\cdot\dot{\vb{r}}) = 2\dot{\vb{r}}\cdot\ddot{\vb{r}} = 0\]
    即只要$\ddot{\vb{r}}\neq\bm{0}$，则$\ddot{\vb{r}}$与$\dot{\vb{r}}$是正交的，故而称$\ddot{\vb{r}}$为曲线$L$的\overtext{主法向量}{principal normal vector}.记：
    \[\vb{\kappa} \coloneq \dot{\vb{r}}\times\ddot{\vb{r}}\]
    显然$\vb{\kappa}$与$\dot{\vb{r}}$也是正交的，称其为$L$的\overtext{副法向量}{binormal vector}.
\end{definition}
设切向量在曲线$L$上一段长度为$\Delta s$弧上的转角为$\Delta\alpha$.在$\Delta\alpha\to0$时，有：
\[\Delta\alpha \sim \tan\Delta\alpha \sim \Delta\dot{\vb{r}} \quad\Leftrightarrow\quad \lim_{\Delta s \to 0}\frac{\Delta\dot{\vb{r}}}{\Delta\alpha} = 1\]
模仿\cref{def:2DCurvature}\textbf{(平面曲率的定义)}，不妨认为空间曲线的曲率也可以这样定义：
\[\kappa = \lim_{\Delta s \to 0}\left|\frac{\Delta\alpha}{\Delta s}\right| = \lim_{\Delta s \to 0}\left|\frac{\Delta\dot{\vb{r}}}{\Delta s}\right| = |\ddot{\vb{r}}| = |\vb{\kappa}|\]
不难推出一般参数表示下曲线的曲率表达式：
\[\vb{\kappa}(t) = \frac{\vb{r}'(t)\times\vb{r}''(t)}{|\vb{r}'(t)|^3}\]
特别地，对于平面曲线，其副法向量始终与$z$轴平行.


\subsection{参数曲面}
参数曲面以向量值函数表示为：
\[\vb{r}(u,v) = x(u,v)\vb{i} + y(u,v)\vb{j} + z(u,v)\vb{k},\quad (u,v)\in D \subset\mathbb{R}^2\]
\subsubsection{切平面}
除了有限个点和曲线段上，设$\vb{r}(u,v)$有连续的偏导数：
\[\begin{cases}
    \displaystyle\vb{r}'_u = \pdv{\vb{r}}{u} = \pdv{x}{u}\vb{i} + \pdv{y}{u}\vb{j} + \pdv{z}{u}\vb{k} \\[1em]
    \displaystyle\vb{r}'_v = \pdv{\vb{r}}{v} = \pdv{x}{v}\vb{i} + \pdv{y}{v}\vb{j} + \pdv{z}{v}\vb{k} \\
\end{cases}\]
显然，根据偏导数的几何性质，$\vb{r}'_u$和$\vb{r}'_v$分别是参数曲面上两不同曲线交点处的两个切向量，而过这一点处的切平面的法向量为$\vb{r}'_u \times \vb{r}'_v$.记单位法向量为：
\[\vb{n} = \frac{\vb{r}'_u \times \vb{r}'_v}{|\vb{r}'_u \times \vb{r}'_v|}\]


\subsubsection{法向量场}
曲面上有一个连续变化的法向量场：
\[\vb{r}'_u(u,v) \times \vb{r}'_v(u,v) = \pdv{(y,z)}{(u,v)}\vb{i} + \pdv{(z,x)}{(u,v)}\vb{j} + \pdv{(x,y)}{(u,v)}\vb{k}\]
这个法向量场中每一点处的法向量的长度实际上是曲面在一点处的切平面上以$\vb{r}'_u$和$\vb{r}'_v$为边的平行四边形的面积，同时也是这一点附近曲面面积的近似.
\begin{definition}
    若参数曲面$\vb{r}(u,v)$有连续的偏导数，且$\vb{r}'_u(u,v) \times \vb{r}'_v(u,v) \neq \bm{0}$，即任意一点附近曲面面积不为$0$.则称这样的曲面是\overtext{光滑曲面}{smooth surface}或\overtext{分片光滑曲面}{piecewise smooth surface}
\end{definition}


\subsection{隐式曲线和隐式曲面}
\subsubsection{平面隐式曲线}
设$F(x,y) \in C^2(D)$，由方程$F(x,y) = 0$可以给出一条隐式曲线.若有$\grad F = F'_x\vb{i} + F'_y\vb{j} \neq \bm{0}$，不妨设有隐函数$y = f(x)$，同时有$f'(x) = -\frac{F'_x(x,y)}{F'_y(x,y)}$.曲线上一点$(x_0,y_0)$的切线和曲率可表示为：
\begin{gather*}
    (x-x_0)f'(x) - (y-y_0) = 0 \quad\Rightarrow\quad (x-x_0)F'_x + (y-y_0)F'_y = 0 \\
    \kappa = \frac{f''(x)}{\left[1+{f'}^2(x)\right]^{3/2}} = \frac{-{F'_y}^2 F''_{xx} + 2F'_x F'_y F''_{xy} - {F'_x}^2 F''_{yy}}{\left[{F'_x}^2 + {F'_y}^2\right]^{3/2}}
\end{gather*}

\begin{mistake}
    上式中的$f''(x)$计算时要注意，$F'_x(x,f(x))$和$F'_y(x,f(x))$都是关于$x$的复合函数，在对其求导时要使用链式法则：
    \[\dv{x}F'_x(x,f(x)) = F''_{xx} + F''_{xy}f'(x),\qquad \dv{x}F'_y(x,f(x)) = F''_{yx} + F''_{yy}f'(x)\]
\end{mistake}

若隐函数以$x = g(y)$表示，得到的结果与前述无异，这表明对于切线方程、曲率，$x,y$是对称的.

从几何上看，方程$F(x,y) = 0$表示了一条等值线，而$F(x,y) = c$所表示的任意等值线由于导数相同，其所具有的切线方程、曲率的表达式与$F(x,y) = 0$时的情形一致.

另一方面，从几何角度观察$\grad F = F'_x\vb{i} + F'_y\vb{j} \neq \bm{0}$这个条件，它可以被重新表述为：只有在有起伏(梯度)时，才存在等值线.


\subsubsection{空间隐式曲面}
设$F(x,y,z) \in C^2(V)$，方程$F(x,y,z) = 0$给出了空间中的一张隐式曲面.从几何上看，$F(x,y,z) = c$是空间中的一张等值面.

如果有一点$(x_0,y_0,z_0)$的$\grad F \neq \bm{0}$，则可以充分地推出在这一点附近有一个连续可导的隐函数，同时曲面上过这一点存在任意的光滑曲线：
\[\vb{r}(t) = x(t)\vb{i} + y(t)\vb{j} + z(t)\vb{k},\quad F(x(t),y(t),z(t)) = 0\]
对$F(x(t),y(t),z(t)) = 0$两端对$t$求导，则可以得到：
\[F'_x x'(t) + F'_y y'(t) + F'_z z'(t) = 0 \quad\Leftrightarrow\quad \grad{F}\cdot\vb{r}'(t) = 0\]
即表明$\grad F$是隐曲面一点处的法向量.


\subsubsection{空间隐式曲线}
空间隐式曲线由两个相交的空间隐式曲面给出.设$F(x,y,z) \in C^2(V),G(x,y,z) \in C^2(V)$，且有$\grad F \neq \bm{0},\grad G \neq \bm{0}$.

曲线一点的切向量$\vb{l}$必然垂直于两个曲面在这一点的法向量，于是有：
\[\vb{l} = \grad F \times \grad G = \pdv{(F,G)}{(y,z)}\vb{i} + \pdv{(F,G)}{(z,x)}\vb{j} + \pdv{(F,G)}{(x,y)}\vb{k}\]


\section{多变量函数的Taylor公式与极值}
\subsection{二元函数的微分中值定理}
\begin{definition}
    称一个区域为\overtext{凸区域}{}，当这个区域中的任意两点的连线仍在这个区域中.
\end{definition}
设凸区域中两点$(x,y)$和$(x+h,y+k)$的线段为$L$：
\[L: \{(x+th,y+tk)\,|\,t\in[0,1]\}\]
用一个函数来表示，就是：
\[\varphi(t) = f(x+th,y+tk),\quad t\in[0,1] \begin{cases}
    \varphi(0) = f(x,y) \\
    \varphi(1) = f(x+h,y+k)
\end{cases}\]
根据一元函数的Lagrange中值公式：
\[\varphi(1) - \varphi(0) = \varphi'(0 + \theta\cdot 1)\cdot 1 = \varphi'(\theta),\quad 0<\theta<1\]
而$\varphi'(t) = hf'_x(x+th,y+tk) + kf'_y(x+th,y+tk)$，由此得到二元函数的中值公式：
\begin{theorem}
    \textbf{二元函数的微分中值公式}

    设函数$f(x,y)$在一个凸区域$D$内有连续的偏导数，则对$D$内任何两点$(x_1,y_1),(x_2,y_2)$存在位于连接它们的直线段上一点$(\xi,\zeta)$使得：
    \[f(x_2,y_2) - f(x_1,y_1) = (x_2-x_1)f'_x(\xi,\zeta) + (y_2-y_1)f'_y(\xi,\zeta)\]
\end{theorem}

由此得到以下推论：
\begin{theorem}
    若函数$f(x,y)$在凸区域$D$内有偏导数且$f'_x = f'_y \equiv 0$，则该函数一定是一个常值函数.

    事实上，只要$D$是道路连通的开集，同样适用上述结论.可以将非凸区域切分成凸区域，连续应用上述结论.
\end{theorem}

模仿上述推导过程，将中值定理推广到$n$维：
\[f(x_1+h_1,\cdots,x_n+h_n) - f(x_1,\cdots,x_n) = \sum_{i=1}^{n}f'_{x_i}(x_1+\theta h_1,\cdots,x_n+\theta h_n)h_i\]

对于向量值函数，微分中值公式不适用，例如：
\[\vb{r}(t) = \cos t\vb{i} + \sin t\vb{j}\]
如果有$\vb{r}(\pi)-\vb{r}(0) = \vb{r}'(\theta)(\pi - 0)$，则会推导出$\sin\theta=\dfrac{2}{\pi}$而$\cos\theta = 0$，这显然不可能.

但是，可以提出一个类似的中值公式：
\begin{theorem}
    \textbf{拟微分中值公式}

    设有$\vb{r}:[a,b]\to\mathbb{R}^2$，有$\vb{r}(t) = x(t)\vb{i} + y(t)\vb{j}$连续可导，则$\exists\theta\in[a,b]$使得：
    \[|\vb{r}(b) - \vb{r}(a)| \leqslant |\vb{r}'(\theta)|(b-a)\]
\end{theorem}

\begin{proof}
    设$\varphi(t) = (\vb{r}(b) - \vb{r}(a))\cdot\vb{r}(t),t\in[a,b]$.有：
    \[\varphi(b) - \varphi(a) = (\vb{r}(b) - \vb{r}(a))\cdot(\vb{r}(b) - \vb{r}(a)) = |\vb{r}(b) - \vb{r}(a)|^2\]
    同时根据一元函数的中值公式，$\exists\theta\in(0,1)$使得：
    \[\varphi(b) - \varphi(a) = \varphi'(\theta)(b-a) = (\vb{r}(b)-\vb{r}(a))\cdot\vb{r}'(\theta)(b-a)\]
    根据\cref{thm:CauchySchwarzInequality}\textbf{(Cauchy-Schwarz不等式)}有：
    \[|\vb{r}(b)-\vb{r}(a)|^2 = |(\vb{r}(b) - \vb{r}(a))\cdot\vb{r}'(\theta)|(b-a) \leqslant |\vb{r}(b)-\vb{r}(a)||\vb{r}'(\theta)|(b-a)\]
    由此得证.
\end{proof}


\subsection{二元函数的Taylor公式}
设$(x_0,y_0)$是区域$D$的内点，在它的凸邻域内函数$f(x,y)\in C^{n+1}$.设有连接$(x_0,y_0)$和$(x,y)\begin{cases}
    x = x_0 + th \\
    y = y_0 + tk
\end{cases}$两点的直线段$\varphi(t) = f(x_0+th,y_0+tk),t\in[0,1]$，求导得：
\[\varphi'(t) = hf'_x(x_0+th,y_0+tk) + kf'_y(x_0+th,y_0+tk) = \mathcal{D}f(x_0+th,y_0+tk)\]
其中微分算子$\displaystyle\mathcal{D}\coloneq h\pdv{x} + k\pdv{y}$.对于二阶导：
\begin{align*}
    \varphi''(t) &= h^2f''_{xx}(x_0+th,y_0+tk) + 2hkf''_{yx}(x_0+th,y_0+tk)+k^2f''_{yy}(x_0+th,y_0+tk) \\
    &= \mathcal{D}^2f(x_0+th,y_0+tk)
\end{align*}
利用数学归纳法不难证明：
\[\varphi^{(m)}(t) = \sum_{i+j=m}\frac{m!}{i!j!}h^i k^j\frac{\partial^m f}{\partial x^i \partial y^j}(x_0+th,y_0+tk) = \mathcal{D}^m f(x_0+th,y_0+tk)\]
对一元函数$\varphi(t)$在$t_0 = 0$处展开，即在$(x,y)$附近展开，得到Maclaurin公式：
\[\varphi(t) = \sum_{m=0}^{n}\frac{\varphi^{(m)}(0)}{m!}t^m + \frac{\varphi^{(n+1)}(\theta t)}{(n+1)!}t^{n+1},\quad 0< \theta < 1\]
取$t = 1$，则有$f(x,y)$在$(x_0,y_0)$处的展开式：
\begin{theorem}
    \textbf{二元函数的Taylor公式}
    \[f(x,y) = \sum_{m=0}^{n}\frac{1}{m!}\mathcal{D}^m f(x_0,y_0) + R_n\]

    其中余项$R_n = \dfrac{1}{(n+1)!}\mathcal{D}^{n+1}f(x_0+\theta h,y_0+\theta k),0<\theta<1$.
\end{theorem}

特别地，由于$f(x,y) \in C^{n+1}$，连续必有界，因此在$(x_0,y_0)$邻域内：
\[\left|\frac{\partial^{n+1}f(x,y)}{\partial x^i \partial y^j}\right| \leqslant M\]
当$h \to 0,k \to 0$时，对$R_n$中$\mathcal{D}^{n+1}$所展开的每一个偏导项都有：
\[\left|\frac{m!}{i!j!}h^i k^j \frac{\partial^m f(x,y)}{\partial x^i \partial y^j}\right| = \frac{m!}{i!j!}|h^i||k^j|\left|\frac{\partial^mf(x,y)}{\partial x^i \partial y^j}\right| \leqslant \frac{m!}{i!j!}|h^i||k^j|M\]
对上式累加求和：
\[|\mathcal{D}^{n+1}f(x,y)| = \left|\sum_{m=0}^{n+1}\frac{m!}{i!j!}h^i k^j \frac{\partial^m f(x,y)}{\partial x^i \partial y^j}\right| \leqslant \sum_{m=0}^{n+1}\frac{m!}{i!j!}|h^i||k^j|M = M(|h|+|k|)^{n+1}\]
进而根据$R_n = \dfrac{\mathcal{D}^{n+1}f(x,y)}{(n+1)!}$：
\[|R_n| \leqslant \frac{M}{(n+1)!}(|h|+|k|)^{n+1}\]
取$\rho = \sqrt{h^2 + k^2}$，有$|h|\leqslant\rho,|k|\leqslant\rho$，于是：
\[|R_n| \leqslant \frac{M}{(n+1)!}(|h|+|k|)^{n+1} \leqslant \frac{2^{n+1}M}{(n+1)!}\rho^{n+1} \quad\Rightarrow\quad 0 \leqslant \frac{|R_n|}{\rho^n} \leqslant \frac{2^{n+1}M}{(n+1)!}\rho\]
在$h\to0,k\to0$时$\rho \to 0$，根据三明治定理：
\[\frac{|R_n|}{\rho^n} \to 0 \quad\Rightarrow\quad \lim_{\rho\to0}\frac{|R_n|}{\rho^n} = 0 \quad\Rightarrow\quad R_n = o(\rho^n)\]
上述Taylor公式的推导及证明过程完全可以推广到$n$维.

特别地，注意到在$(X_1,\cdots,X_n)$处$n$维一阶Taylor展开的公式为：
\begin{align*}
    f(x_1,\cdots,x_n) &= f(X_1,\cdots,X_n) + \sum_{i=1}^{n}h_i f'_{x_i} + R_1 \\
    &= f(X_1,\cdots,X_n) + (h_1,\cdots,h_n)\cdot\grad{f} + R_1
\end{align*}
定义$\vb{h} \coloneq (h_1,\cdots,h_n)$，称$\vb{H}_f \coloneq \begin{bmatrix}
    f''_{x_1 x_1} & \cdots & f''_{x_1 x_n} \\
    \vdots & \ddots & \vdots \\
    f''_{x_n x_1} & \cdots & f''_{x_n x_n}
\end{bmatrix}$为\textbf{Hessian矩阵}，上式可写为如下二次型的形式：
\[f(x_1,\cdots,x_n) = f(X_1,\cdots,X_n) + \vb{h}\vb{H}_f\vb{h}^{\top} + R_1\]


\subsection{二元函数的极值}
\begin{definition}
    设$f(x)$在区域$D$中有定义，$M_0(x_0,y_0) \in D$.若在$M_0$的邻域内恒有：
    \[f(x,y) \leqslant f(x_0,y_0) \quad\text{或}\quad f(x,y) \geqslant f(x_0,y_0)\]
    则称$f(x_0,y_0)$是$f(x,y)$的一个\textbf{极大值}或\textbf{极小值}，$M_0$称为极值点.
\end{definition}

\begin{theorem}
    若$f(x,y)$有偏导数，$(x_0,y_0)$是极值点，则有$f'_x(x_0,y_0) = f'_y(x_0,y_0) = 0$.反之不然.
\end{theorem}

\begin{definition}
    在一点使函数一阶偏导数为$0$的点称为函数的\textbf{驻点}，是成为极值点的必要条件.
\end{definition}